\documentstyle [11pt]{article}
\pagestyle{plain}
\bibliographystyle{plain}

\input amssym.def
\input amssym.tex
%\def\baselinestretch{2}\small\normalsize

\setlength{\textwidth}{152.4truemm}
\setlength{\textheight}{228.6truemm}
\setlength{\oddsidemargin}{3.6mm}
\setlength{\evensidemargin}{3.6mm}
\setlength{\topmargin}{-12.5truemm}
\setlength{\parindent}{5.0truemm}

%\newtheorem{theorem}{Theorem}[section]
\newtheorem{theorem}{Theorem}
\newtheorem{lemma}[theorem]{Lemma}
\newtheorem{corollary}[theorem]{Corollary}
\newtheorem{example}[theorem]{Example}
\newtheorem{definition}[theorem]{Definition}
\newtheorem{remark}[theorem]{Remark}

\def\arraystretch{1.153}                 %You can change 1.0
\def\n#1{\vbox to 3mm{\vspace{1mm}\vfill \hbox to 2.0mm{\hfill
             $#1$\hfill} \vfill }}


\title{A census of critical sets in the Latin squares \\
of order at most six}
\author{Peter Adams, Richard Bean and Abdollah Khodkar \\
Centre for Discrete Mathematics and Computing \\ 
Department of Mathematics\\ 
The University of Queensland\\ 
Queensland 4072\\ Australia} 

\date{ }


\begin{document}
\maketitle


\vspace{15 pt}

\begin{quote}
        ABSTRACT: A critical set in a Latin square of order $n$ 
	is a set of entries in a Latin square which can be embedded 
	in precisely one Latin square of order $n$. Also, if any element 
	of the critical set is deleted, the remaining set can be 
	embedded in more than one Latin square of order $n$.
	In this paper we find all the critical sets of different 
	sizes in the Latin squares of order at most six. 
        We count the number of main and isotopy classes of these critical
        sets and classify critical sets from the main classes into various
        ``strengths''.  Some observations are made about the relationship
        between the numbers of classes, particularly in the $6 \times 6$
        case.  Finally some examples are given of each type of critical set.
\end{quote}

\footnotetext[1]{Research supported by the Australian Research Council}
%%%%%%%%%%%%%%%%%%%%%%%%%%%%%%%%%%%SECTION1
\section{Introduction}   
\setcounter{theorem}{0}
A {\it Latin square} $L$ of order $n$ is an 
$n\times n$ array with entries 
chosen from a set $N$, of size $n$, such that each element of $N$ occurs 
precisely once in each row and column. 
For convenience, a Latin square will sometimes be represented as a 
set of ordered triples $(i,j;k)$.  This means that element 
% changed from "this is read to mean" to "this means"; sentence split
$k$ occurs in cell $(i,j)$ of the Latin square $L$.
For a Latin square of order $n$ we shall use $1,2,3,\dots,n$ 
as the entries, and rows and columns will also be labelled from 
$1,\dots,n$.
A {\it partial Latin square} $P$ of order $n$ is an $n\times n$ 
array with entries chosen from a set $N$, of size $n$, such that each
element of $N$ occurs at most once in each row and column. Thus $P$ may
contain a number of empty cells. We sometimes denote $P$ by 
the set $\{(i,j;k)|i,j,k\in N\}$.
If a Latin square $L$ contains an $s\times s$ 
subarray $S$ and if $S$ is a Latin square of order $s$, then we say 
that $S$ is a {\it Latin subsquare} of $L$. 
A {\it critical set} in a Latin square $L$ of order $n$, is a 
set ${\cal C}=\{(i,j;k)\;|\; i,j,k\in N\}$, such that
\begin{itemize}
\item[(1)] $L$ is the only Latin square of order $n$ which has 
element $k$ in cell $(i,j)$ for each $(i,j;k)\in {\cal C}$; and

\item[(2)] no proper subset of ${\cal C}$ satisfies $(1)$.
\end{itemize}

A partial Latin square $U$ of $L$ which satisfies condition (1) is called
a {\em uniquely completable} (UC) set in $L$.  A {\it smallest critical set} 
in a Latin square $L$ is a critical set of minimum cardinality.

Let $L$ be a Latin square of order $n$ and let
$\{a,b,c\}=\{1,2,3\}$. The $(a,b,c)$-{\it conjugate} of $L$, written
$L_{(a,b,c)}$, is defined as follows:
\[ L_{(a,b,c)}=\{(x_a,x_b;x_c)\mid (x_1,x_2;x_3)\in L\}. \]
Two Latin squares $L$ and $L'$ of order $n$ are 
{\it isotopic} if there are three bijections from the rows, columns, and
symbols of $L$ to the rows, columns, and symbols, respectively, of 
$L'$, that map $L$ to $L'$. Two Latin squares $L$ and $L'$ of 
order $n$ are {\it main class isotopic} if $L$ is isotopic to any 
conjugate of $L'$. Table 1 shows the number of main and isotopy
classes for Latin squares of order $1\leq n\leq 7$ (see Denes and
Keedwell \cite {dk}). 

\begin{table}
\begin{center}
\caption{Number of main and isotopy classes for Latin squares of order 1 to 7}
\begin{tabular}{|c|c|c|c|c|c|c|c|} \hline
$n=$ & $1$ & $2$ & $3$ & $4$ & $5$ & $6$ & $7$ \\ \hline
Main classes & $1$ & $1$ & $1$ & $2$ & $2$ & $12$ & $147$ \\ \hline
Isotopic classes & $1$ & $1$ & $1$ & $2$ & $2$ & $22$ & $564$ \\ \hline
\end{tabular} \vspace {3mm} \\
\end{center}
\end{table}

In the process of completing the critical set ${\cal C}$ to the 
Latin square $L$ of order $n$ which it characterizes, we say that adjunction 
of a triple $t=(r,c;s)$ is {\em forced} (see \cite {MR99j:05036})
in the process of completion of a set $T$ of triples 
$(|T|<n^2,\;{\cal C}\subseteq T\subset L)$ to the complete set of
triples which represents $L$, if either

\begin{itemize}
\item [(i)] $\forall r' \neq r$, $\exists z \neq c$ such that 
$(r',z;s)\in T$ or $\exists z \neq s$ such that 
$(r',c;z)\in T$, or 

\item [(ii)] $\forall c' \neq c$, $\exists z \neq r$ such that 
$(z,c';s)\in T$ or $\exists z \neq s$ such that 
$(r,c';z)\in T$, or

\item [(iii)] $\forall s' \neq s$, $\exists z \neq r$ such that 
$(z,c;s')\in T$ or $\exists z \neq c$ such that 
$(r,z;s')\in T$.
\end{itemize}

\noindent The critical set ${\cal C}$ is called {\em strong} 
if we can define a sequence of sets of triples 
${\cal C} = F_1 \subset F_2\subset F_3\subset \dots \subset F_r=L$ 
such that each triple $t\in F_{i+1}\setminus F_i$ is forced in 
$F_i$ for $1\leq i\leq r-1$. 
%A critical set is {\em super-strong} if each triple in this 
%sequence is forced only by virtue of 
%property (iii)in the above definition of forcing. 
A critical set which is not strong is called {\em weak}.
A critical set is called {\em totally weak} if no cell is forced.
A {\it smallest weak critical set} in a Latin square $L$ is 
a weak critical set of minimum cardinality.
(Note that the authors of \cite {MR2000g:05034} have used 
the term {\em semi-strong} in place of strong 
as defined above.)

The following extension of the concept of the strong critical set is
taken from Bedford and Whitehouse ~\cite{bandw}.

Let $P$ be a partial Latin square of order $n$ defined on an element set $N$.  
Then $A_P$ is an {\it array of alternatives} for $P$ if

\begin{enumerate}
\item $A_P$ is an $n \times n$ array;
\item whenever the cell at $(i,j)$ in $P$ is filled, the cell at $(i,j)$ of $A_P$ is empty; and
\item whenever the cell at $(i,j)$ in $P$ is empty, the cell at $(i,j)$ of $A_P$ contains all the elements of $N$ which do not appear in the $i^{th}$ row or $j^{th}$ column of $P$.
\end{enumerate}
%We denote the entry in row $i$ and column $j$ of $A_P$ by $(i,j)_{A_P}$. 
We denote the set of elements in cell $(i,j)$ of $A_P$ by ${A_P}(i,j)$.

Let $P$ be a partial Latin square.  
%The addition of a triple $t = (i,j;k)$
%is said to be {\it semi-forced} if for each $k' \neq k$ in $(i,j)_{A_P}$, either
%
%\begin{enumerate}
%\item there exists $r > 0$ and $i_1, \dots, i_r$ (all $\neq i$) with $k' \in 
%(i_1,j)_{A_P} \cup \dots \cup (i_r,j)_{A_P}$ and $\mid (i_1,j)_{A_P} \cup \dots
%\cup (i_r,j)_{A_P} \mid = r$, or
%\item there exists $s > 0$ and $j_1, \dots, j_s$ (all $\neq j$) with $k' \in
%(i,j_1)_{A_P} \cup \dots \cup (i,j_s)_{A_P}$ and $\mid (i,j_1)_{A_P} \cup \dots \cup (i,j_s)_{A_P} \mid = s$.
%\end{enumerate}
We shall say that the symbol
$k'\in A_P(i,j)$ is {\it forced out} of $A_P$ if either:
\begin{itemize}
\item [(1)] there exists $r>0$ and $i_1,i_2,\ldots,i_r$ (all $\neq i$)
with $k'\in A_P(i_1,j)\cup\ldots\cup A_P(i_r,j)$ and
$|A_P(i_1,j)\cup\ldots\cup A_P(i_r,j)|=r$; or
\item [(2)] $\theta(i,j,k')$ satisfies $1$ in $A_{P_{\theta(1,2,3)}}$ for
some $\theta\in S_3$.
\end{itemize}

The reduced array of alternatives,
$RA_P$, is the array obtained from $A_P$ by successively removing
symbols which are forced out until no more symbols can be forced out.
Then the addition
of an entry $(i,j;k)$ to $P$ is said to be {\it semi-forced} if either:

\begin{enumerate}
\item $k$ is the only symbol in $RA_P(i,j)$; or
\item $k$ occurs exactly once in either the $i^{\rm {th}}$ row or
$j^{\rm {th}}$ column of $RA_P$.
\end{enumerate}

Then a UC set $U$ is {\it near-strong} to the Latin square $L$
if we can find a sequence of sets $U = S_1 \subset \dots \subset S_f = L$
such that each triple $t \in S_{v+1} \setminus S_{v}$ is semi-forced
in $S_v$.  We call a UC set {\it Bedford-Whitehouse totally weak} if no
cell is semi-forced.  If a UC set is Bedford-Whitehouse totally weak,
this implies that it is also totally weak.

We shall denote the number of critical sets of size $x$ in a main class
$y.z$ by $CS(y,z,x)$. (The notation $y.z$ denotes main class $z$ in
a Latin square of order $y$, as in the CRC Handbook of Combinatorial
Designs,~\cite{crc}.)  The number of isotopy classes of critical sets
of size $x$ in a main class $y.z$ shall be denoted $IC(y,z,x)$, and
the number of main classes of these critical sets shall be denoted
$MC(y,z,x)$.  The greatest common divisor of the number of critical
sets of all sizes in a particular main class $y.z$ will be referred to
as GCDCS$(y,z)$.

The sizes of smallest critical sets for the Latin squares of orders 
four and five were determined in \cite {MR95k:05030,MR80j:05022}.
Howse in \cite {MR98m:05026} finds 
smallest critical sets for all the Latin squares of order six.
Adams and Khodkar in \cite {AK2} give smallest critical sets for 
all the Latin squares of order at most seven. They also find \cite {MR2000k:05054}
smallest weak and smallest totally weak critical sets for the Latin squares 
of order at most seven.
The size of smallest strong 
critical sets in a Latin square 
has also been considered in the past (see \cite {MR2000g:05034}). 

This paper deals with critical sets of different sizes in the Latin
squares of order at most six.  First, for each main class of Latin
square of order at most six, we calculate every possible critical set.
(These will be of various sizes.)  Then, again for each main class and
possible size, the main classes and isotopy classes for the relevant set
of critical sets are determined.  Next, we determine which of the main
classes are strong, near-strong, totally weak, and
Bedford-Whitehouse totally weak.  Some interesting properties concerning
the greatest common divisors of numbers of critical sets across main classes in
the $6 \times 6$ Latin squares and ratios of various kinds are discussed.
%
%Finally, for each main class
%of Latin square examined, the possibility of the Latin square being a
%union of critical sets for the same Latin square is examined.

In the appendix, we list examples for each possible size of critical set in
each main class of $6 \times 6$ Latin square.

\vspace{5mm}

%%%%%%%%%%%%%%%%%%%%%%%%%%%%%SECTION algorithm
\section{Algorithms}

There were two basic algorithms used to independently calculate all critical sets
of a given size $m$ for a given main class of $n \times n$ Latin squares.  
The first examined all ${n^2 \choose m}$
% used {n^2 \choose m} instead of \binom{n^2}{m} with amstex
size $m$ subsets of the
Latin square.  If any subset $U$ had unique completion, all the proper
subsets of $U$ which are of size $\mid U \mid -1$ were tested for
unique completion.  If no such subset had unique completion, $U$ was output
as a critical set.

The second algorithm used Latin interchanges (see ~\cite{MR98b:05019})
to speed up this process.  Interchanges are subsets of a Latin square
which have the property that any critical set must intersect every 
Latin interchange in at least one cell.  This program divided the Latin square
up into disjoint Latin interchanges, ensuring that each candidate for a
critical set had at least one element in each of the Latin interchanges.

For the $6 \times 6$ case, for critical sets of size greater than 18,
the process was further speeded up by ensuring that in
each subset examined, no row or column was full and no symbol 
occurred six times.  Such subsets cannot be critical sets as any
cell may be removed from the relevant row, column or symbol set
while maintaining the unique completion property.


%%%%%%%%%%%%%%%%%%%%%%%%%%%%%SECTION2
\section{Table of results}

\subsection{Explanation of headings}
The first column in the tables of results (Tables \ref{one}, \ref{two}, 
\ref{three} and \ref{four}) is
the main class number $y.z$ (LS), followed by the size(s) of the critical
set(s) for the main class (Size), the number of critical sets of that
size in the main class (\#CS); this is then followed by the number of
isotopy classes (\#Iso) and the number of main classes of those critical sets (\#Main).
(The notation $y.z$ denotes main class $z$ in
a Latin square of order $y$, as in the CRC Handbook of Combinatorial
Designs,~\cite{crc}.)  

%The last four columns list the number of critical sets of various ``strengths''
%within the main classes of critical sets.  

For the next four columns, we consider representatives of each main
class of critical sets, and list the number of critical sets of various
``strengths'' within the main classes of critical sets.  That is,
we calculate how many of the representatives of each main class of
critical set have the various ``strengths''.  We need only consider
representatives of each main class of critical set, since, for example,
a strong critical set remains a strong critical set when the rows, columns and
symbols are interchanged or swapped.  Similarly, a 
near-strong critical set remains near-strong  under
permutations or interchanges of rows, columns or symbols.
These last four columns are, in order, the number of 
near-strong critical sets (\#NS), the number of strong critical sets
(\#Strong), the number of totally weak critical sets (\#TW), and the
number of Bedford-Whitehouse totally weak critical sets (\#BWTW).
\subsection{Latin squares of order 3}

There is only one main class, denoted 3.1, for Latin squares of order three \cite{crc}: 
\vspace{5mm}

\begin{center}
\begin{tabular}{|c|c|c|} \hline
1&2&3 \\ \hline
2&3&1 \\ \hline
3&1&2 \\ \hline
\end{tabular} \vspace{3mm}\\
{\bf 3.1}
\end{center}
\vspace{5mm}

For this Latin square, we have the results presented in Table \ref{one} 
concerning the number of critical sets of every possible size.

\begin{table}
\begin{center}
\caption{Critical set statistics for Latin squares of order 3}
\label{one}
\begin{tabular}{|c|c|c|c|c|c|c|c|c|} \hline
LS & Size & \#CS & \#Iso & \#Main & \#NS & \#Strong & \#TW & \#BWTW \\ \hline
$3.1$  & $2$ & $9$ & $1$ & $1$ & 1 & 1 & 0 & 0  \\
& $3$ & 18 & $1$ & $1$ & 1 & 1 & 0 & 0 \\ \hline
\end{tabular}\vspace {3mm} \\
\end{center}
\end{table}

\subsection{Latin squares of order 4}

There are two main classes, denoted 4.1 and 4.2, for Latin squares of order four \cite{crc}:
\vspace{5mm}

\begin{center}
\begin{tabular}{cccc}
  \begin{tabular}{|c|c|c|c|} \hline
  1&2&3&4 \\ \hline
  2&3&4&1 \\ \hline
  3&4&1&2 \\ \hline
  4&1&2&3 \\ \hline
  \end{tabular} &&&
  \begin{tabular}{|c|c|c|c|} \hline
  1&2&3&4 \\ \hline
  2&1&4&3  \\ \hline
  3&4&1&2  \\ \hline
  4&3&2&1 \\ \hline
  \end{tabular} \\
  {\bf 4.1} &&& {\bf 4.2} \\
\end{tabular}
\end{center}
\vspace{5mm}

For these Latin squares, we have the results presented in Table \ref{two} 
concerning the number of critical sets of every possible size.

\begin{table}
\begin{center}
\caption{Critical set statistics for Latin squares of order 4}
\label{two}
\begin{tabular}{|c|c|c|c|c|c|c|c|c|} \hline
LS & Size & \#CS & \#Iso & \#Main & \#NS & \#Strong & \#TW & \#BWTW \\ \hline
4.1 & $4$ & $32$  & $1$  & $1$ & 1 & 1 & 0 & 0\\
    & $5$ & $576$ & $18$ & $4$ & 4 & 4 & 0 & 0\\
    & $6$ & $128$ & $4$  & $2$ & 2 & 2 & 0 & 0\\ \hline
4.2 & $5$ & $96$  & $1$  & $1$ & 1 & 1 & 0 & 0\\
    & $6$ & $432$ & $7$  & $3$ & 3 & 3 & 0 & 0\\
    & $7$ & $48$  & $1$ & $1 $ & 1 & 1 & 0 & 0\\ \hline
\end{tabular}\vspace {3mm} \\
\end{center}
\end{table}

\subsection{Latin squares of order 5}

There are two main classes, denoted 5.1 and 5.2, for Latin squares of order five \cite{crc}:

\vspace{5mm}

\begin{center}
\begin{tabular}{cccc}
  \begin{tabular}{|c|c|c|c|c|} \hline
  1&2&3&4&5  \\ \hline
  2&3&4&5&1  \\ \hline
  3&4&5&1&2  \\ \hline
  4&5&1&2&3  \\ \hline
  5&1&2&3&4   \\ \hline
  \end{tabular} &&&
  \begin{tabular}{|c|c|c|c|c|} \hline
  1&2&3&4&5 \\ \hline
  2&1&4&5&3 \\ \hline
  3&4&5&1&2 \\ \hline
  4&5&2&3&1 \\ \hline
  5&3&1&2&4 \\ \hline
  \end{tabular} \\
  {\bf 5.1} &&& {\bf 5.2} \\
\end{tabular}
\end{center}
\vspace{5mm}

\begin{table}
\begin{center}
\caption{Critical set statistics for Latin squares of order 5}
\label{three}
\begin{tabular}{|c|c|c|c|c|c|c|c|c|} \hline
 LS & Size & \#CS &
\#Iso & \#Main & \#NS & \#Strong & \#TW & \#BWTW \\ \hline
$5.1$ & $6$  & $50$    & $1$    & $1$ & 1 & 1 & 0 & 0 \\
      & $7$  & $1000$  & $10$   & $4$ & 4 & 4  & 0 & 0 \\
      & $8$  & $30900$ & $312$  & $57$ & 57 & 57 & 0 & 0 \\
      & $9$  & $18800$ & $188$  & $37$ & 37 & 37 & 0 & 0 \\
      & $10$ & $2500$  & $25$   & $6$ & 6 & 6 & 0 & 0 \\ \hline
$5.2$ & $7$  & $600$   & $50$   & $11$ & 10 & 10 & 1 & 1 \\
      & $8$  & $21588$ & $1802$ & $322$ & 311 & 311 & 1 & 1 \\
      & $9$  & $23718$ & $1981$ & $348$ & 348 & 348 & 0 & 0 \\
      & $10$ & $2340$  & $198$  & $39$ & 38 & 36 & 2 & 0 \\
      & $11$ & $216$   & $18$   & $4$ & 4 & 4 & 0 & 0 \\ \hline
\end{tabular}\vspace {3mm} \\
\end{center}
\end{table}

\subsection{Latin squares of order 6}

There are twelve main classes, denoted 6.1, \dots, 6.12, for Latin squares of 
order six \cite{crc}:

\begin{center}
\begin{tabular}{cccc}
\begin{tabular}{|c@{\hspace{1.5mm}}|c@{\hspace{1.5mm}}|c@{\hspace{1.5mm}}
|c@{\hspace{1.5mm}}|c@{\hspace{1.5mm}}|c@{\hspace{1.5mm}}|} \hline
1&2&3&4&5&6  \\ \hline
2&3&4&5&6&1  \\ \hline
3&4&5&6&1&2  \\ \hline
4&5&6&1&2&3  \\ \hline
5&6&1&2&3&4  \\ \hline
6&1&2&3&4&5  \\ \hline
\end{tabular}&
\begin{tabular}{|c@{\hspace{1.5mm}}|c@{\hspace{1.5mm}}|c@{\hspace{1.5mm}}
|c@{\hspace{1.5mm}}|c@{\hspace{1.5mm}}|c@{\hspace{1.5mm}}|} \hline
1&2&3&4&5&6  \\ \hline
2&1&4&3&6&5  \\ \hline
3&4&5&6&1&2  \\ \hline
4&3&6&5&2&1  \\ \hline
5&6&1&2&4&3  \\ \hline
6&5&2&1&3&4  \\ \hline
\end{tabular}&
\begin{tabular}{|c@{\hspace{1.5mm}}|c@{\hspace{1.5mm}}|c@{\hspace{1.5mm}}
|c@{\hspace{1.5mm}}|c@{\hspace{1.5mm}}|c@{\hspace{1.5mm}}|} \hline
1&2&3&4&5&6  \\ \hline
2&1&4&5&6&3  \\ \hline
3&4&1&6&2&5  \\ \hline
4&5&6&1&3&2  \\ \hline
5&6&2&3&1&4  \\ \hline
6&3&5&2&4&1  \\ \hline
\end{tabular} &
\begin{tabular}{|c@{\hspace{1.5mm}}|c@{\hspace{1.5mm}}|c@{\hspace{1.5mm}}
|c@{\hspace{1.5mm}}|c@{\hspace{1.5mm}}|c@{\hspace{1.5mm}}|} \hline
1&2&3&4&5&6  \\ \hline
2&1&4&5&6&3  \\ \hline
3&4&1&6&2&5  \\ \hline
4&5&6&1&3&2  \\ \hline
5&6&2&3&4&1  \\ \hline
6&3&5&2&1&4  \\ \hline
\end{tabular} \\ 
{\bf 6.1} & {\bf 6.2} & {\bf 6.3} & {\bf 6.4}\\
\end{tabular}
\end{center}

\begin{center}
\begin{tabular}{cccc}
\begin{tabular}{|c@{\hspace{1.5mm}}|c@{\hspace{1.5mm}}|c@{\hspace{1.5mm}}
|c@{\hspace{1.5mm}}|c@{\hspace{1.5mm}}|c@{\hspace{1.5mm}}|} \hline
1&2&3&4&5&6  \\ \hline
2&1&4&5&6&3  \\ \hline
3&4&2&6&1&5  \\ \hline
4&5&6&2&3&1  \\ \hline
5&6&1&3&4&2  \\ \hline
6&3&5&1&2&4  \\ \hline
\end{tabular} & 
\begin{tabular}{|c@{\hspace{1.5mm}}|c@{\hspace{1.5mm}}|c@{\hspace{1.5mm}}
|c@{\hspace{1.5mm}}|c@{\hspace{1.5mm}}|c@{\hspace{1.5mm}}|} \hline
1&2&3&4&5&6  \\ \hline
2&1&4&5&6&3  \\ \hline
3&4&5&6&1&2  \\ \hline
4&5&6&3&2&1  \\ \hline
5&6&1&2&3&4  \\ \hline
6&3&2&1&4&5  \\ \hline
\end{tabular} & 
\begin{tabular}{|c@{\hspace{1.5mm}}|c@{\hspace{1.5mm}}|c@{\hspace{1.5mm}}
|c@{\hspace{1.5mm}}|c@{\hspace{1.5mm}}|c@{\hspace{1.5mm}}|} \hline
1&2&3&4&5&6  \\ \hline
2&1&4&3&6&5  \\ \hline
3&5&1&6&2&4  \\ \hline
4&6&2&5&1&3  \\ \hline
5&3&6&1&4&2  \\ \hline
6&4&5&2&3&1  \\ \hline
\end{tabular} & 
\begin{tabular}{|c@{\hspace{1.5mm}}|c@{\hspace{1.5mm}}|c@{\hspace{1.5mm}}
|c@{\hspace{1.5mm}}|c@{\hspace{1.5mm}}|c@{\hspace{1.5mm}}|} \hline
1&2&3&4&5&6  \\ \hline
2&1&4&3&6&5  \\ \hline
3&5&1&6&2&4  \\ \hline
4&6&2&5&1&3  \\ \hline
5&3&6&2&4&1  \\ \hline
6&4&5&1&3&2  \\ \hline
\end{tabular} \\
{\bf 6.5} & {\bf 6.6} & {\bf 6.7} & {\bf 6.8}\\
\end{tabular}
\end{center}

\begin{center}
\begin{tabular}{cccc}
\begin{tabular}{|c@{\hspace{1.5mm}}|c@{\hspace{1.5mm}}|c@{\hspace{1.5mm}}
|c@{\hspace{1.5mm}}|c@{\hspace{1.5mm}}|c@{\hspace{1.5mm}}|} \hline
1&2&3&4&5&6  \\ \hline
2&1&4&3&6&5  \\ \hline
3&5&1&6&2&4  \\ \hline
4&6&2&5&3&1  \\ \hline
5&4&6&2&1&3  \\ \hline
6&3&5&1&4&2  \\ \hline
\end{tabular} & 
\begin{tabular}{|c@{\hspace{1.5mm}}|c@{\hspace{1.5mm}}|c@{\hspace{1.5mm}}
|c@{\hspace{1.5mm}}|c@{\hspace{1.5mm}}|c@{\hspace{1.5mm}}|} \hline
1&2&3&4&5&6  \\ \hline
2&1&4&3&6&5  \\ \hline
3&5&1&6&4&2  \\ \hline
4&6&5&1&2&3  \\ \hline
5&3&6&2&1&4  \\ \hline
6&4&2&5&3&1  \\ \hline
\end{tabular} & 
\begin{tabular}{|c@{\hspace{1.5mm}}|c@{\hspace{1.5mm}}|c@{\hspace{1.5mm}}
|c@{\hspace{1.5mm}}|c@{\hspace{1.5mm}}|c@{\hspace{1.5mm}}|} \hline
1&2&3&4&5&6  \\ \hline
2&1&4&5&6&3  \\ \hline
3&4&2&6&1&5  \\ \hline
4&6&5&2&3&1  \\ \hline
5&3&6&1&2&4  \\ \hline
6&5&1&3&4&2  \\ \hline
\end{tabular} & 
\begin{tabular}{|c@{\hspace{1.5mm}}|c@{\hspace{1.5mm}}|c@{\hspace{1.5mm}}
|c@{\hspace{1.5mm}}|c@{\hspace{1.5mm}}|c@{\hspace{1.5mm}}|} \hline
1&2&3&4&5&6  \\ \hline
2&3&1&5&6&4  \\ \hline
3&1&2&6&4&5  \\ \hline
4&6&5&2&1&3  \\ \hline
5&4&6&3&2&1  \\ \hline
6&5&4&1&3&2  \\ \hline
\end{tabular} \\
{\bf 6.9} & {\bf 6.10} & {\bf 6.11} & {\bf 6.12}\\
\end{tabular}
\end{center}
%\vspace{5mm}

For these Latin squares, we have the results presented in Table \ref{four} 
concerning the number
of critical sets of every possible size.

\begin{table}
\begin{center}
\caption{Critical set statistics for Latin squares of order 6}
\label{four}
\begin{tabular}{|c|c|c|c|c|c|c|c|c|} \hline
 LS & Size & \#CS & \#Iso & \#Main& \#NS & \#Strong & \#TW & \#BWTW \\ \hline
$6.1$ & $9$  & $72$    & $1$    & $1$ & 1 & 1 & 0 & 0 \\
      & $11$  & $39384$ & $547$  & $97$ & 97 & 95 & 0 & 0  \\
      & $12$  & $1161036$ & $16149$  & $2740$ & 2541 & 2513 & 11 & 8 \\
      & $13$  & $3634344$ & $50492$  & $8481$ & 7815 & 7792 & 19 & 14 \\
      & $14$  & $886428$ & $12346$  & $2090$ & 1931 & 1920 & 10 & 9 \\
      & $15$  & $80064$ & $1118$  & $202$ & 182 & 168 & 8 & 0 \\
      & $16$  & $3240$ & $45$  & $8$ & 8 & 8 & 0 & 0 \\
      & $17$  & $108$  & $3$   & $1$ & 0 & 0 & 0 & 0 \\ \hline
$6.2$ & $11$  & $7848$ & $327$  & $60$ & 50 & 48 & 3 & 3 \\
      & $12$  & $658908$ & $27477$  & $4633$ & 4370 & 4325 & 35 & 27 \\
      & $13$  & $3328908$ & $138708$  & $23267$ & 22226 & 22187 & 52 & 36 \\
      & $14$  & $1800228$ & $75035$  & $12617$ & 12267 & 12263 & 11 & 8 \\
      & $15$  & $192480$ & $8022$  & $1362$ & 1354 & 1351 & 3 & 1 \\
      & $16$  & $15840$ & $660$  & $115$ & 113 & 113 & 0 & 0 \\
      & $17$  & $240$  & $10$   & $3$ & 3 & 3 & 0 & 0  \\ \hline
$6.3$ & $11$  & $1200$ & $10$  & $7$ & 2 & 2 & 0 & 0 \\
      & $12$  & $192360$ & $1603$  & $836$ & 749 & 748 & 14 & 14 \\
      & $13$  & $1837440$ & $15315$  & $7757$ & 7445 & 7440 & 33 & 29 \\
      & $14$  & $1727880$ & $14400$  & $7279$ & 7252 & 7252 & 2 & 2 \\
      & $15$  & $378928$ & $3162$  & $1610$ & 1610 & 1610 & 0 & 0 \\
      & $16$  & $20280$ & $169$  & $90$ & 90 & 90 & 0 & 0 \\
      & $17$  & $840$  & $7$   & $4$ & 4 & 4 & 0 & 0 \\ \hline
\end{tabular}\vspace {3mm} \\
\end{center}
\end{table}
\newpage
\begin{center}
Table \ref{four} (continued)
\begin{tabular}{|c|c|c|c|c|c|c|c|c|} \hline
 LS & Size & \#CS & \#Iso & \#Main& \#NS & \#Strong & \#TW & \#BWTW \\ \hline
$6.4$ & $10$  & $56$ & $7$  & $5$ & 5 & 5 & 0 & 0 \\
      & $11$  & $34000$ & $4250$  & $2149$ & 2001 & 1980 & 16 & 10 \\
      & $12$  & $1590608$ & $198826$  & $99654$ & 94485 & 94024 & 197 & 136 \\
      & $13$  & $5498076$ & $687262$  & $344044$ & 328754 & 327801 & 331 & 232 \\
      & $14$  & $1931424$ & $241428$  & $120895$ & 116390 & 115691 & 58 & 34 \\
      & $15$  & $168752$ & $21095$  & $10586$ & 10102 & 9704 & 137 & 10 \\
      & $16$  & $13736$ & $1717$  & $871$ & 821 & 780 & 24 & 1 \\
      & $17$  & $148$  & $19$   & $11$ & 9 & 9 & 1 & 1 \\ \hline
$6.5$ & $10$  & $60$ & $15$  & $9$ & 9 & 9 & 0 & 0 \\
      & $11$  & $42992$ & $10748$  & $5406$ & 5132 & 5078 & 30 & 24 \\
      & $12$  & $1878236$ & $469559$  & $235063$ & 224705 & 223776 & 401 & 292 \\
      & $13$  & $6475142$ & $1618806$  & $809952$ & 778258 & 776251 & 648 & 473 \\
      & $14$  & $2182652$ & $545663$  & $273120$ & 264790 & 263229 & 119 & 76 \\
      & $15$  & $192416$ & $48104$  & $24108$ & 23304 & 22340 & 281 & 28 \\
      & $16$  & $16908$ & $4227$  & $2135$ & 2041 & 1961 & 43 & 3 \\
      & $17$  & $112$  & $28$   & $16$ & 15 & 15 & 0 & 0 \\ \hline
$6.6$ & $11$  & $12888$ & $358$  & $187$ & 177 & 175 & 0 & 0 \\
      & $12$  & $856908$ & $23803$  & $12005$ & 11191 & 11155 & 64 & 55 \\
      & $13$  & $4097790$ & $113839$  & $57151$ & 54038 & 53898 & 162 & 120 \\
      & $14$  & $1476864$ & $41024$  & $20664$ & 19770 & 19697 & 27 & 23 \\
      & $15$  & $155196$ & $4311$  & $2201$ & 2139 & 2117 & 8 & 4 \\
      & $16$  & $12744$ & $354$  & $186$ & 175 & 166 & 3 & 0 \\
      & $17$  & $216$  & $6$   & $4$ & 4 & 4 & 0 & 0 \\ \hline
\end{tabular} \vspace{3mm} \\
\end{center}
\newpage
\begin{center}
Table \ref{four} (continued)
\begin{tabular}{|c|c|c|c|c|c|c|c|c|} \hline
LS & Size & \#CS & \#Iso & \#Main& \#NS & \#Strong & \#TW & \#BWTW \\ \hline
$6.7$ & $12$  & $4752$ & $22$ & $5$ & 3 & 3 & 0 & 0 \\
      & $13$  & $212328$ & $985$ & $183$ & 165 & 165 & 5 & 5 \\
      & $14$  & $893700$ & $4151$ & $736$ & 706 & 705 & 3 & 2 \\
      & $15$  & $545508$ & $2536$ & $465$ & 465 & 465 & 0 & 0 \\
      & $16$  & $125766$ & $583$ & $109$ & 109 & 109 & 0 & 0 \\
      & $17$  & $8208$ & $38$ & $13$ & 13 & 13 & 0 & 0 \\ 
      & $18$  & $648$ & $3$ & $1$ & 1 & 1 & 0 & 0 \\ \hline
$6.8$ & $11$  & $3264$ & $408$ & $75$ & 67 & 66 & 3 & 2 \\
      & $12$  & $324608$ & $40576$ & $6817$ & 6023 & 5986 & 37 & 33 \\
      & $13$  & $1826592$ & $228335$ & $38265$ & 35161 & 35063 & 123 & 80 \\
      & $14$  & $1093796$ & $136729$ & $22909$ & 21804 & 21764 & 10 & 8 \\
      & $15$  & $106296$ & $13290$ & $2260$ & 2178 & 2155 & 10 & 1 \\
      & $16$  & $8464$ & $1058$ & $185$ & 175 & 167 & 6 & 1 \\ 
      & $17$  & $216$ & $27$ & $5$ & 5 & 5 & 0 & 0 \\ \hline
$6.9$ & $10$  & $24$ & $2$ & $2$ & 2 & 2 & 0 & 0 \\
      & $11$  & $13980$ & $1165$ & $596$ & 546 & 535 & 7 & 7 \\
      & $12$  & $716352$ & $59714$ & $29939$ & 27999 & 27723 & 127 & 84 \\
      & $13$  & $2784264$ & $232027$ & $116246$ & 109378 & 108885 & 243 & 173 \\
      & $14$  & $1065876$ & $88856$ & $44575$ & 42345 & 42068 & 24 & 19 \\
      & $15$  & $85884$ & $7159$ & $3607$ & 3462 & 3314 & 72 & 4 \\
      & $16$  & $6960$ & $580$ & $302$ & 283 & 259 & 15 & 1 \\ 
      & $17$  & $24$ & $2$ & $2$ & 2 & 2 & 0 & 0 \\ \hline
\end{tabular}\vspace {3mm} \\
\end{center}
\newpage
\begin{center}
Table \ref{four} (continued)
\begin{tabular}{|c|c|c|c|c|c|c|c|c|} \hline
 LS & Size & \#CS & \#Iso & \#Main& \#NS & \#Strong & \#TW & \#BWTW \\ \hline
$6.10$& $10$  & $4$ & $1$ & $1$ & 1 & 1 & 0 & 0 \\
      & $11$  & $13748$ & $3437$ & $587$ & 555 & 547 & 6 & 3 \\
      & $12$  & $858348$ & $214587$ & $35899$ & 33814 & 33644 & 169 & 123 \\
      & $13$  & $3894038$ & $973520$ & $162538$ & 154803 & 154404 & 279 & 195 \\
      & $14$  & $1715492$ & $428873$ & $71685$ & 69560 & 69375 & 38 & 29 \\
      & $15$  & $155000$ & $38753$ & $6513$ & 6355 & 6232 & 34 & 4 \\
      & $16$  & $10540$ & $2635$ & $461$ & 443 & 423 & 13 & 1 \\ 
      & $17$  & $120$ & $30$ & $6$ & 6 & 6 & 0 & 0 \\ \hline
$6.11$& $10$  & $40$ & $10$ & $3$ & 3 & 3 & 0 & 0 \\
      & $11$  & $63540$ & $15885$ & $2673$ & 2617 & 2590 & 9 & 5 \\
      & $12$  & $2292266$ & $573254$ & $95781$ & 92453 & 92029 & 96 & 59 \\
      & $13$  & $7075888$ & $1768972$ & $295196$ & 284917 & 284027 & 145 & 96 \\
      & $14$  & $2203696$ & $550993$ & $91977$ & 88209 & 87499 & 39 & 22 \\
      & $15$  & $188344$ & $47086$ & $7890$ & 7584 & 7175 & 161 & 8  \\
      & $16$  & $17172$ & $4293$ & $729$ & 685 & 645 & 35 & 4 \\ 
      & $17$  & $36$ & $9$ & $2$ & 1 & 1 & 0 & 0 \\ \hline
$6.12$& $11$  & $143208$ & $1326$ & $232$ & 229 & 228 & 0 & 0 \\
      & $12$  & $3518478$ & $32664$ & $5510$ & 5384 & 5358 & 7 & 6 \\
      & $13$  & $9025344$ & $83568$ & $14037$ & 13636 & 13584 & 17 & 14 \\
      & $14$  & $2104704$ & $19506$ & $3315$ & 3146 & 3107 & 6 & 4 \\
      & $15$  & $200340$ & $1855$ & $326$ & 316 & 297 & 8 & 1 \\
      & $16$  & $17820$ & $165$ & $32$ & 29 & 28 & 1 & 0 \\ \hline
\end{tabular}\vspace {3mm} \\
\end{center}

%\newpage
% so that the caption "Table 3" is associated with the table
% keedwell, pages 124-128
%Keedwell ~\cite{dk} points out that (for isotopy and main classes of Latin squares)
%each isotopy class has a number of Latin squares
%associated with it, and similarly each main class has a number of isotopy
%classes associated with it.  This also applies for main classes and
%isotopy classes of critical sets.  Each main class of critical sets
%in $6 \times 6$ Latin squares has either 1, 2, 3, or 6 associated
%isotopy classes.  Each isotopy class of critical sets in $6 \times 6$
%Latin squares has anywhere from 2 to 216 associated critical sets.
%These results are given in Tables 6 and 7 below.  

D{\'e}nes and Keedwell ~\cite{dk} point out that, for a given order $n$, each
isotopy class of $n \times n$ Latin squares has a number of Latin
squares associated with it, and similarly each main class of $n \times
n$ Latin squares has a number of isotopy classes associated with it.
Similarly, for any given main class $n.z$ of $n \times n$ Latin squares
and given size of critical set $m$, if we consider the main classes of
critical sets of size $m$ within the main class $n.z$, we have several
associated isotopy classes of critical sets.  In the same way, if we
consider the isotopy classes of critical sets of size $m$ within the
main class $n.z$, we have several associated critical sets of size $m$
in the main class $n.z$.
%\newpage

In Tables \ref{five} to \ref{eight}, the head line refers to the twelve main cla
sses of $6 \times
6$ Latin squares, and the side line refers to the possible sizes of
critical sets.  For $6 \times 6$ Latin squares, we consider results related to t
hese observations.

Each isotopy class of critical sets in $6 \times 6$ Latin squares has between
2 to 216 associated critical sets.  This result is given in Table \ref{five}.


\begin{table}
\begin{center}
\begin{footnotesize}
\caption{Numbers of critical sets in each isotopy class of critical sets of order six}
\label{five}
\begin{tabular}{|c@{\hspace{1.0mm}}|c@{\hspace{1.0mm}}|c@{\hspace{1.0mm}}|c@{\hspace{1.0mm}}|c@{\hspace{1.0mm}}|c@{\hspace{1.0mm}}|c@{\hspace{1.0mm}}|c@{\hspace{1.0mm}}|c@{\hspace{1.0mm}}|c@{\hspace{1.0mm}}|c@{\hspace{1.0mm}}|c@{\hspace{1.0mm}}|c@{\hspace{1.0mm}}|} \hline
& $6.1$ & $6.2$ & $6.3$ & $6.4$ & $6.5$ & $6.6$ & $6.7$ & $6.8$ & $6.9$ & $6.10$ & $6.11$ & $6.12$ \\
\hline
9 &  72 &   -   &    -   &  -   &  -   &   -   &    -   &  -   & -     &  -   &  -   &  - \\
10 &   -   &   -   &    -   & 8 & 4 &   -   &    -   &  -   & 12 & 4 & 4 &  - \\
11 & 72 & 24 & 120 & 8 & 4 & 36 &    -   & 8 & 12 & 4 & 4 & 108  \\
12 & 12,36,72 & 8,12,24 & 8  & 8 & 4 & 36 & 216 & 8 & 6,12 & 4 & 2,4 & 54,108 \\
13 & 36,72 & 12,24 & 60,120 & 4,8 & 2,4 & 18,36 & 108,216 & 4,8 & 6,12 & 2,4 & 4 & 108 \\
14 & 36,72 & 12,24 & 60,120 & 8 & 4  & 36 & 108,216 & 4,8 & 6,12 & 4 & 2,4 & 54,108 \\
15 & 36,72 & 12,24 & 24,40,120 & 4,8 & 4 & 36 & 108,216 & 4,8 & 6,12 & 2,4 & 4 & 108 \\
16 & 72 & 24 & 120 & 8 & 4 & 36 & 54,216 & 8 & 12 & 4 & 4 & 108 \\
17 & 36 & 24 & 120 & 4,8 & 4 & 36 & 216 & 8 & 12 & 4 & 4 &    - \\
18 &   -   &   -   &    -   &  -   &  -   &   -   & 216 &  -   &   -   &  -   &  -   &    -  \\
\hline
\end{tabular}\vspace {3mm} \\
\end{footnotesize}
\end{center}
\end{table}

Each main class of critical sets in $6 \times 6$ Latin squares has either
1, 2, 3 or 6 associated isotopy classes of critical sets.  
This result is given in Table \ref{six}. 

\begin{table}
\begin{center}
\caption{Numbers of isotopy classes of critical sets in each main class of critical sets of order six \vspace{3mm}}
\label{six}
\begin{tabular}{|c@{\hspace{1.0mm}}|c@{\hspace{1.0mm}}|c@{\hspace{1.0mm}}|c@{\hspace{1.0mm}}|c@{\hspace{1.0mm}}|c@{\hspace{1.0mm}}|c@{\hspace{1.0mm}}|c@{\hspace{1.0mm}}|c@{\hspace{1.0mm}}|c@{\hspace{1.0mm}}|c@{\hspace{1.0mm}}|c@{\hspace{1.0mm}}|c@{\hspace{1.0mm}}|} \hline
& $6.1$ & $6.2$ & $6.3$ & $6.4$ & $6.5$ & $6.6$ & $6.7$ & $6.8$ & $6.9$ & $6.10$ & $6.11$ & $6.12$ \\
\hline
9 &  1 &   -   &    -   &  -   &  -   &   -   &    -   &  -   & -     &  -   &  -   &  - \\
10 &   -   &   -   &    -   & 1,2 & 1,2 &   -   &    -   &  -   & 1 & 1 & 1,3,6 &  - \\
11 & 1,3,6 & 3,6 & 1,2 & 1,2 & 1,2 & 1,2 &    -   & 3,6 & 1,2  & 2,3,6 & 1,2,3,6 & 3,6  \\
12 & 1,2,3,6 & 1,2,3,6 & 1,2 & 1,2 & 1,2 & 1,2 & 1,3,6 & 1,2,3,6 & 1,2 & 1,2,3,6 & 1,2,3,6 & 1,2,3,6 \\
13 & 2,3,6 & 1,2,3,6 & 1,2 & 1,2 & 1,2 & 1,2 & 1,3,6 & 1,2,3,6 & 1,2 & 1,2,3,6 & 1,2,3,6 & 3,6 \\
14 & 1,2,3,6 & 1,2,3,6 & 1,2 & 1,2 & 1,2 & 1,2 & 1,2,3,6 & 1,2,3,6 & 1,2 & 2,3,6 & 2,3,6 & 3,6 \\
15 & 1,2,3,6 & 1,2,3,6 & 1,2 & 1,2 & 1,2 & 1,2 & 1,3,6 & 3,6 & 1,2 & 1,2,3,6 & 1,3,6 & 1,2,3,6 \\
16 & 3,6 & 3,6 & 1,2 & 1,2 & 1,2 & 1,2 & 1,3,6 & 2,3,6 & 1,2 & 1,2,3,6 & 3,6 & 3,6 \\
17 & 3 & 1,3,6 & 1,2 & 1,2 & 1,2 & 1,2 & 1,3,6 & 3,6 & 1 & 3,6 & 3,6 &    - \\
18 &   -   &   -   &    -   &  -   &  -   &   -   & 3 &  -   &   -   &  -   &  -   &    -  \\
\hline
\end{tabular}\vspace {3mm} \\
\end{center}
\end{table}
 
\section{Some observations}

We will concentrate on observations concerning the $6 \times 6$ Latin squares.

We find that when the main class $z$ is fixed and $x$ takes all possible
values, $CS(6,z,x)$ / $IC(6,z,x)$ is in most cases close to an integer
constant.  There is one exception: the critical sets of size 17 in main
class 6.1, where all other values of $CS(6,1,x) / IC(6,1,x)$ are approximately 72,
but $CS(6,1,17) / IC(6,1,17) = 36$.  We also find that this integer constant is a
multiple of GCDCS$(6,z)$.  These ratios are given in Table \ref{seven} above,
truncated at two decimal places.  The bottom line tabulates the values
of GCDCS$(6,z)$.

\begin{table}
\begin{center}
\begin{footnotesize}
\caption{Ratio of critical sets to isotopy classes of critical sets of order six}
\label{seven}
\begin{tabular}{|c|c|c|c|c|c|c|c|c|c|c|c|c|} \hline
& $6.1$ & $6.2$ & $6.3$ & $6.4$ & $6.5$ & $6.6$ & $6.7$ & $6.8$ & $6.9$ & $6.10$ & $6.11$ & $6.12$ \\
\hline
9 &  72.00 &   -   &    -   &  -   &  -   &   -   &    -   &  -   & -     &  -   &  -   &  - \\
10 &   -   &   -   &    -   & 8.00 & 4.00 &   -   &    -   &  -   & 12.00 & 4.00 & 4.00 &  - \\
11 & 72.00 & 24.00 & 120.00 & 8.00 & 4.00 & 36.00 &    -   & 8.00 & 12.00 & 4.00 & 4.00 & 108.00  \\
12 & 71.89 & 23.98 & 120.00 & 8.00 & 4.00 & 36.00 & 216.00 & 8.00 & 11.99 & 4.00 & 3.99 & 107.71 \\
13 & 71.97 & 23.99 & 119.97 & 7.99 & 3.99 & 35.99 & 215.56 & 7.99 & 11.99 & 3.99 & 4.00 & 108.00 \\
14 & 71.79 & 23.99 & 119.99 & 8.00 & 4.00 & 36.00 & 215.29 & 7.99 & 11.99 & 4.00 & 3.99 & 107.90 \\
15 & 71.61 & 23.99 & 119.83 & 7.99 & 4.00 & 36.00 & 215.10 & 7.99 & 11.99 & 3.99 & 4.00 & 108.00 \\
16 & 72.00 & 24.00 & 120.00 & 8.00 & 4.00 & 36.00 & 215.72 & 8.00 & 12.00 & 4.00 & 4.00 & 108.00 \\
17 & 36.00 & 24.00 & 120.00 & 7.78 & 4.00 & 36.00 & 216.00 & 8.00 & 12.00 & 4.00 & 4.00 &    - \\
18 &   -   &   -   &    -   &  -   &  -   &   -   & 216.00 &  -   &   -   &  -   &  -   &    -  \\
\hline
gcd & 36 & 12 & 2 & 4 & 2 & 18 & 54 & 4 & 12 & 2 & 2 & 54  \\
\hline
\end{tabular}\vspace {3mm} \\
\end{footnotesize}
\end{center}
\end{table}

In seven of the twelve $6 \times 6$ main classes (6.1, 6.2, 6.7, 6.8,
6.10, 6.11, and 6.12), the ratio $MC(6,z,x)$/$IC(6,z,x)$ is close to 6 with a few exceptions.
In the other five main classes (6.3, 6.4, 6.5, 6.6, and 6.9) this ratio
is close to 2 with a few exceptions.  These ratios are given in Table \ref{eight} above, truncated at
two decimal places.

\begin{table}
\begin{center}
\caption{Ratio of main classes of critical sets to isotopy classes of critical sets of order six \vspace {3mm}}
\label{eight}
\begin{tabular}{|c|c|c|c|c|c|c|c|c|c|c|c|c|} \hline
& $6.1$ & $6.2$ & $6.3$ & $6.4$ & $6.5$ & $6.6$ & $6.7$ & $6.8$ & $6.9$ & $6.10$ & $6.11$ & $6.12$ \\
\hline
9 &   1.00 &   -  &  -   &  -   &  -   &  -   &  -   &  -   &   -  &  -   &  -   &  - \\
10 &   -   &   -  &  -   & 1.40 & 1.66 &  -   &  -   &  -   & 1.00 & 1.00 & 3.33 &  - \\
11 &  5.63 & 5.45 & 1.42 & 1.97 & 1.98 & 1.91 &  -   & 5.44 & 1.95 & 5.85 & 5.94 & 5.71  \\
12 &  5.89 & 5.93 & 1.91 & 1.99 & 1.99 & 1.98 & 4.40 & 5.95 & 1.99 & 5.97 & 5.98 & 5.92 \\
13 &  5.95 & 5.96 & 1.97 & 1.99 & 1.99 & 1.99 & 5.38 & 5.96 & 1.99 & 5.98 & 5.99 & 5.95 \\
14 &  5.90 & 5.94 & 1.97 & 1.99 & 1.99 & 1.98 & 5.63 & 5.96 & 1.99 & 5.98 & 5.99 & 5.88 \\
15 &  5.53 & 5.88 & 1.96 & 1.99 & 1.99 & 1.95 & 5.45 & 5.88 & 1.98 & 5.95 & 5.96 & 5.69 \\
16 &  5.62 & 5.73 & 1.87 & 1.97 & 1.97 & 1.90 & 5.34 & 5.71 & 1.92 & 5.71 & 5.88 & 5.15 \\
17 &  3.00 & 3.33 & 1.75 & 1.72 & 1.75 & 1.50 & 2.92 & 5.40 & 1.00 & 5.00 & 4.50 &  - \\
18 &   -   &   -  &  -   &  -   &  -   &  -   & 3.00 &  -   &  -   &  -   &  -   &  -  \\
\hline
\end{tabular}\vspace {3mm} \\
\end{center}
\end{table}

In each $6 \times 6$ main class $6.z$, GCDCS$(6,z)$ is a multiple of
2, and for those main classes with $3 \times 3$ subsquares (6.1, 6.6,
6.7 and 6.12) this number is a multiple of 18.

We also note that in the $4 \times 4$ and $6 \times 6$ main classes of
Latin squares, the smallest and largest possible critical sets (4 and 7
for the $4 \times 4$ case and 9 and 18 for the $6 \times 6$ case) each
have only one isotopy and main class.

\section{Appendix}
In this appendix, we give examples of critical sets of each possible size 
in each of the twelve main classes of $6 \times 6$ Latin squares.
\begin{center}
\begin{tabular}{c@{\hspace{1.0mm}}c@{\hspace{1.0mm}}c@{\hspace{1.0mm}}c@{\hspace{1.0mm}}}
\begin{tabular}{|c|c|c|c|c|c|}
\hline 1 & 2 & 3 & $\;\;$ & $\;\;$ & $\;\;$ \\
\hline 2 & 3 & $\;\;$ & $\;\;$ & $\;\;$ & $\;\;$ \\
\hline 3 & $\;\;$ & $\;\;$ & $\;\;$ & $\;\;$ & $\;\;$ \\
\hline $\;\;$ & $\;\;$ & $\;\;$ & $\;\;$ & $\;\;$ & $\;\;$ \\
\hline $\;\;$ & $\;\;$ & $\;\;$ & $\;\;$ & $\;\;$ & 4 \\
\hline $\;\;$ & $\;\;$ & $\;\;$ & $\;\;$ & 4 & 5 \\
\hline 
\end{tabular} &
\begin{tabular}{|c|c|c|c|c|c|}
\hline $\;\;$ & $\;\;$ & $\;\;$ & $\;\;$ & $\;\;$ & $\;\;$ \\
\hline $\;\;$ & $\;\;$ & $\;\;$ & $\;\;$ & $\;\;$ & 1 \\
\hline $\;\;$ & $\;\;$ & $\;\;$ & $\;\;$ & 1 & 2 \\
\hline $\;\;$ & $\;\;$ & $\;\;$ & 1 & 2 & 3 \\
\hline $\;\;$ & $\;\;$ & 1 & 2 & 3 & 4 \\
\hline 6 & $\;\;$ & $\;\;$ & $\;\;$ & $\;\;$ & $\;\;$ \\
\hline 
\end{tabular} &
\begin{tabular}{|c|c|c|c|c|c|}
\hline $\;\;$ & $\;\;$ & $\;\;$ & $\;\;$ & $\;\;$ & $\;\;$ \\
\hline $\;\;$ & $\;\;$ & $\;\;$ & $\;\;$ & $\;\;$ & 1 \\
\hline $\;\;$ & $\;\;$ & $\;\;$ & $\;\;$ & 1 & 2 \\
\hline $\;\;$ & $\;\;$ & $\;\;$ & 1 & 2 & 3 \\
\hline $\;\;$ & 6 & 1 & 2 & 3 & $\;\;$ \\
\hline 6 & $\;\;$ & $\;\;$ & $\;\;$ & $\;\;$ & 5 \\
\hline 
\end{tabular} &
\begin{tabular}{|c|c|c|c|c|c|}
\hline $\;\;$ & $\;\;$ & $\;\;$ & $\;\;$ & $\;\;$ & $\;\;$ \\
\hline $\;\;$ & $\;\;$ & $\;\;$ & $\;\;$ & $\;\;$ & 1 \\
\hline $\;\;$ & $\;\;$ & $\;\;$ & $\;\;$ & 1 & 2 \\
\hline $\;\;$ & $\;\;$ & 6 & 1 & 2 & $\;\;$ \\
\hline $\;\;$ & 6 & 1 & 2 & 3 & $\;\;$ \\
\hline $\;\;$ & $\;\;$ & $\;\;$ & 3 & 4 & 5 \\
\hline 
\end{tabular}\\ 
6.1, size 9
 & 6.1, size 11
 & 6.1, size 12
 & 6.1, size 13
\\
\end{tabular} 
\end{center} 
\begin{center}
\begin{tabular}{c@{\hspace{1.0mm}}c@{\hspace{1.0mm}}c@{\hspace{1.0mm}}c@{\hspace{1.0mm}}}
\begin{tabular}{|c|c|c|c|c|c|}
\hline $\;\;$ & $\;\;$ & $\;\;$ & $\;\;$ & $\;\;$ & $\;\;$ \\
\hline $\;\;$ & $\;\;$ & 1 & $\;\;$ & $\;\;$ & 4 \\
\hline $\;\;$ & 1 & 2 & $\;\;$ & 4 & 5 \\
\hline $\;\;$ & $\;\;$ & $\;\;$ & $\;\;$ & $\;\;$ & 1 \\
\hline $\;\;$ & $\;\;$ & 4 & 3 & 1 & 2 \\
\hline $\;\;$ & 4 & 5 & $\;\;$ & 2 & $\;\;$ \\
\hline 
\end{tabular} &
\begin{tabular}{|c|c|c|c|c|c|}
\hline $\;\;$ & $\;\;$ & $\;\;$ & $\;\;$ & $\;\;$ & $\;\;$ \\
\hline $\;\;$ & $\;\;$ & $\;\;$ & $\;\;$ & $\;\;$ & 1 \\
\hline $\;\;$ & $\;\;$ & $\;\;$ & $\;\;$ & 1 & 2 \\
\hline $\;\;$ & $\;\;$ & $\;\;$ & 1 & 2 & 3 \\
\hline $\;\;$ & $\;\;$ & 1 & 2 & 3 & 4 \\
\hline $\;\;$ & 1 & 2 & 3 & 4 & 5 \\
\hline 
\end{tabular} &
\begin{tabular}{|c|c|c|c|c|c|}
\hline $\;\;$ & $\;\;$ & $\;\;$ & $\;\;$ & $\;\;$ & $\;\;$ \\
\hline $\;\;$ & $\;\;$ & $\;\;$ & $\;\;$ & $\;\;$ & 1 \\
\hline $\;\;$ & $\;\;$ & 5 & $\;\;$ & 1 & 2 \\
\hline $\;\;$ & 5 & 6 & 1 & 2 & $\;\;$ \\
\hline $\;\;$ & 6 & 1 & 2 & $\;\;$ & 4 \\
\hline $\;\;$ & 1 & 2 & 3 & $\;\;$ & 5 \\
\hline 
\end{tabular} &
\begin{tabular}{|c|c|c|c|c|c|}
\hline $\;\;$ & $\;\;$ & $\;\;$ & $\;\;$ & $\;\;$ & $\;\;$ \\
\hline $\;\;$ & $\;\;$ & $\;\;$ & 5 & 6 & 1 \\
\hline $\;\;$ & $\;\;$ & 5 & $\;\;$ & 1 & 2 \\
\hline $\;\;$ & 5 & 6 & 1 & 2 & 3 \\
\hline $\;\;$ & 6 & 1 & $\;\;$ & 3 & $\;\;$ \\
\hline $\;\;$ & 1 & 2 & 3 & $\;\;$ & $\;\;$ \\
\hline 
\end{tabular}\\ 
6.1, size 14
 & 6.1, size 15
 & 6.1, size 16
 & 6.1, size 17
\\
\end{tabular} 
\end{center} 
\begin{center}
\begin{tabular}{c@{\hspace{1.0mm}}c@{\hspace{1.0mm}}c@{\hspace{1.0mm}}c@{\hspace{1.0mm}}}
\begin{tabular}{|c|c|c|c|c|c|}
\hline $\;\;$ & $\;\;$ & $\;\;$ & $\;\;$ & $\;\;$ & $\;\;$ \\
\hline $\;\;$ & 1 & $\;\;$ & 3 & $\;\;$ & 5 \\
\hline $\;\;$ & $\;\;$ & $\;\;$ & 6 & 1 & $\;\;$ \\
\hline $\;\;$ & 3 & 6 & $\;\;$ & 2 & $\;\;$ \\
\hline 5 & $\;\;$ & $\;\;$ & $\;\;$ & $\;\;$ & $\;\;$ \\
\hline $\;\;$ & $\;\;$ & 2 & $\;\;$ & $\;\;$ & 4 \\
\hline 
\end{tabular} &
\begin{tabular}{|c|c|c|c|c|c|}
\hline $\;\;$ & $\;\;$ & $\;\;$ & $\;\;$ & $\;\;$ & $\;\;$ \\
\hline $\;\;$ & 1 & $\;\;$ & 3 & $\;\;$ & 5 \\
\hline $\;\;$ & $\;\;$ & $\;\;$ & $\;\;$ & 1 & 2 \\
\hline $\;\;$ & 3 & 6 & $\;\;$ & $\;\;$ & $\;\;$ \\
\hline 5 & $\;\;$ & $\;\;$ & 2 & $\;\;$ & $\;\;$ \\
\hline 6 & $\;\;$ & 2 & $\;\;$ & $\;\;$ & 4 \\
\hline 
\end{tabular} &
\begin{tabular}{|c|c|c|c|c|c|}
\hline $\;\;$ & $\;\;$ & $\;\;$ & $\;\;$ & $\;\;$ & $\;\;$ \\
\hline $\;\;$ & 1 & $\;\;$ & 3 & $\;\;$ & 5 \\
\hline $\;\;$ & $\;\;$ & $\;\;$ & $\;\;$ & 1 & 2 \\
\hline $\;\;$ & 3 & $\;\;$ & 5 & 2 & $\;\;$ \\
\hline $\;\;$ & 6 & 1 & $\;\;$ & 4 & $\;\;$ \\
\hline $\;\;$ & $\;\;$ & 2 & $\;\;$ & $\;\;$ & 4 \\
\hline 
\end{tabular} &
\begin{tabular}{|c|c|c|c|c|c|}
\hline $\;\;$ & $\;\;$ & $\;\;$ & $\;\;$ & $\;\;$ & $\;\;$ \\
\hline $\;\;$ & 1 & $\;\;$ & 3 & $\;\;$ & 5 \\
\hline $\;\;$ & $\;\;$ & $\;\;$ & $\;\;$ & 1 & 2 \\
\hline $\;\;$ & 3 & $\;\;$ & 5 & 2 & $\;\;$ \\
\hline $\;\;$ & $\;\;$ & 1 & 2 & 4 & $\;\;$ \\
\hline $\;\;$ & 5 & 2 & $\;\;$ & $\;\;$ & 4 \\
\hline 
\end{tabular}\\ 
6.2, size 11
 & 6.2, size 12
 & 6.2, size 13
 & 6.2, size 14
\\
\end{tabular} 
\end{center} 
\begin{center}
\begin{tabular}{c@{\hspace{1.0mm}}c@{\hspace{1.0mm}}c@{\hspace{1.0mm}}c@{\hspace{1.0mm}}}
\begin{tabular}{|c|c|c|c|c|c|}
\hline $\;\;$ & $\;\;$ & $\;\;$ & $\;\;$ & $\;\;$ & $\;\;$ \\
\hline $\;\;$ & 1 & $\;\;$ & 3 & $\;\;$ & 5 \\
\hline $\;\;$ & $\;\;$ & $\;\;$ & $\;\;$ & 1 & 2 \\
\hline $\;\;$ & 3 & $\;\;$ & 5 & 2 & $\;\;$ \\
\hline $\;\;$ & $\;\;$ & 1 & 2 & 4 & $\;\;$ \\
\hline 6 & $\;\;$ & 2 & 1 & 3 & $\;\;$ \\
\hline 
\end{tabular} &
\begin{tabular}{|c|c|c|c|c|c|}
\hline $\;\;$ & $\;\;$ & $\;\;$ & $\;\;$ & $\;\;$ & $\;\;$ \\
\hline $\;\;$ & 1 & $\;\;$ & 3 & $\;\;$ & 5 \\
\hline $\;\;$ & $\;\;$ & $\;\;$ & $\;\;$ & 1 & 2 \\
\hline $\;\;$ & 3 & 6 & 5 & 2 & 1 \\
\hline $\;\;$ & 6 & 1 & $\;\;$ & $\;\;$ & $\;\;$ \\
\hline $\;\;$ & 5 & 2 & 1 & 3 & $\;\;$ \\
\hline 
\end{tabular} &
\begin{tabular}{|c|c|c|c|c|c|}
\hline $\;\;$ & $\;\;$ & $\;\;$ & $\;\;$ & $\;\;$ & $\;\;$ \\
\hline $\;\;$ & 1 & $\;\;$ & 3 & $\;\;$ & 5 \\
\hline $\;\;$ & $\;\;$ & $\;\;$ & $\;\;$ & 1 & 2 \\
\hline $\;\;$ & 3 & $\;\;$ & 5 & 2 & 1 \\
\hline $\;\;$ & $\;\;$ & 1 & 2 & $\;\;$ & 3 \\
\hline $\;\;$ & 5 & 2 & 1 & 3 & 4 \\
\hline 
\end{tabular} &
\begin{tabular}{|c|c|c|c|c|c|}
\hline $\;\;$ & $\;\;$ & $\;\;$ & $\;\;$ & 5 & 6 \\
\hline $\;\;$ & 1 & $\;\;$ & 3 & $\;\;$ & $\;\;$ \\
\hline $\;\;$ & $\;\;$ & 1 & $\;\;$ & $\;\;$ & $\;\;$ \\
\hline 4 & $\;\;$ & $\;\;$ & $\;\;$ & $\;\;$ & 2 \\
\hline $\;\;$ & 4 & $\;\;$ & $\;\;$ & 1 & $\;\;$ \\
\hline 6 & $\;\;$ & 2 & $\;\;$ & $\;\;$ & $\;\;$ \\
\hline 
\end{tabular}\\ 
6.2, size 15
 & 6.2, size 16
 & 6.2, size 17
 & 6.3, size 11
\\
\end{tabular} 
\end{center} 
\begin{center}
\begin{tabular}{c@{\hspace{1.0mm}}c@{\hspace{1.0mm}}c@{\hspace{1.0mm}}c@{\hspace{1.0mm}}}
\begin{tabular}{|c|c|c|c|c|c|}
\hline $\;\;$ & $\;\;$ & $\;\;$ & $\;\;$ & $\;\;$ & 6 \\
\hline $\;\;$ & 1 & $\;\;$ & 3 & $\;\;$ & $\;\;$ \\
\hline $\;\;$ & $\;\;$ & 1 & $\;\;$ & 2 & $\;\;$ \\
\hline $\;\;$ & $\;\;$ & 5 & 1 & $\;\;$ & $\;\;$ \\
\hline 5 & $\;\;$ & $\;\;$ & 2 & $\;\;$ & $\;\;$ \\
\hline 6 & $\;\;$ & 2 & $\;\;$ & 4 & $\;\;$ \\
\hline 
\end{tabular} &
\begin{tabular}{|c|c|c|c|c|c|}
\hline $\;\;$ & $\;\;$ & $\;\;$ & $\;\;$ & $\;\;$ & 6 \\
\hline $\;\;$ & 1 & $\;\;$ & 3 & $\;\;$ & $\;\;$ \\
\hline $\;\;$ & $\;\;$ & 1 & $\;\;$ & 2 & $\;\;$ \\
\hline $\;\;$ & $\;\;$ & $\;\;$ & 1 & 3 & 2 \\
\hline 5 & 4 & $\;\;$ & $\;\;$ & $\;\;$ & $\;\;$ \\
\hline $\;\;$ & 3 & $\;\;$ & 5 & 4 & $\;\;$ \\
\hline 
\end{tabular} &
\begin{tabular}{|c|c|c|c|c|c|}
\hline $\;\;$ & $\;\;$ & $\;\;$ & $\;\;$ & $\;\;$ & 6 \\
\hline $\;\;$ & 1 & $\;\;$ & 3 & $\;\;$ & $\;\;$ \\
\hline $\;\;$ & $\;\;$ & 1 & $\;\;$ & 2 & $\;\;$ \\
\hline $\;\;$ & $\;\;$ & $\;\;$ & 1 & 3 & 2 \\
\hline $\;\;$ & 4 & $\;\;$ & $\;\;$ & 1 & 3 \\
\hline 6 & $\;\;$ & 2 & $\;\;$ & 4 & $\;\;$ \\
\hline 
\end{tabular} &
\begin{tabular}{|c|c|c|c|c|c|}
\hline $\;\;$ & $\;\;$ & $\;\;$ & $\;\;$ & $\;\;$ & $\;\;$ \\
\hline $\;\;$ & 1 & $\;\;$ & 3 & $\;\;$ & 5 \\
\hline $\;\;$ & $\;\;$ & 1 & $\;\;$ & 2 & 4 \\
\hline $\;\;$ & $\;\;$ & $\;\;$ & 1 & 3 & 2 \\
\hline $\;\;$ & $\;\;$ & 6 & 2 & 1 & $\;\;$ \\
\hline 6 & $\;\;$ & 2 & $\;\;$ & 4 & $\;\;$ \\
\hline 
\end{tabular}\\ 
6.3, size 12
 & 6.3, size 13
 & 6.3, size 14
 & 6.3, size 15
\\
\end{tabular} 
\end{center} 
\begin{center}
\begin{tabular}{c@{\hspace{1.0mm}}c@{\hspace{1.0mm}}c@{\hspace{1.0mm}}c@{\hspace{1.0mm}}}
\begin{tabular}{|c|c|c|c|c|c|}
\hline $\;\;$ & $\;\;$ & $\;\;$ & $\;\;$ & $\;\;$ & $\;\;$ \\
\hline $\;\;$ & 1 & $\;\;$ & 3 & $\;\;$ & 5 \\
\hline $\;\;$ & $\;\;$ & 1 & $\;\;$ & 2 & 4 \\
\hline $\;\;$ & $\;\;$ & $\;\;$ & 1 & 3 & 2 \\
\hline $\;\;$ & $\;\;$ & 6 & 2 & 1 & $\;\;$ \\
\hline $\;\;$ & 3 & $\;\;$ & 5 & 4 & 1 \\
\hline 
\end{tabular} &
\begin{tabular}{|c|c|c|c|c|c|}
\hline $\;\;$ & $\;\;$ & $\;\;$ & $\;\;$ & $\;\;$ & $\;\;$ \\
\hline $\;\;$ & 1 & $\;\;$ & 3 & $\;\;$ & 5 \\
\hline $\;\;$ & $\;\;$ & 1 & $\;\;$ & 2 & 4 \\
\hline $\;\;$ & $\;\;$ & $\;\;$ & 1 & 3 & 2 \\
\hline $\;\;$ & 4 & $\;\;$ & $\;\;$ & 1 & 3 \\
\hline $\;\;$ & 3 & 2 & 5 & 4 & 1 \\
\hline 
\end{tabular} &
\begin{tabular}{|c|c|c|c|c|c|}
\hline $\;\;$ & $\;\;$ & $\;\;$ & 4 & $\;\;$ & 6 \\
\hline $\;\;$ & 1 & $\;\;$ & $\;\;$ & $\;\;$ & $\;\;$ \\
\hline $\;\;$ & 5 & 1 & $\;\;$ & $\;\;$ & $\;\;$ \\
\hline 4 & $\;\;$ & $\;\;$ & $\;\;$ & $\;\;$ & 3 \\
\hline $\;\;$ & $\;\;$ & $\;\;$ & $\;\;$ & $\;\;$ & $\;\;$ \\
\hline 6 & $\;\;$ & 2 & $\;\;$ & $\;\;$ & 4 \\
\hline 
\end{tabular} &
\begin{tabular}{|c|c|c|c|c|c|}
\hline $\;\;$ & $\;\;$ & $\;\;$ & $\;\;$ & $\;\;$ & $\;\;$ \\
\hline $\;\;$ & 1 & $\;\;$ & 3 & $\;\;$ & 5 \\
\hline $\;\;$ & $\;\;$ & 1 & $\;\;$ & $\;\;$ & 2 \\
\hline $\;\;$ & $\;\;$ & $\;\;$ & 1 & $\;\;$ & 3 \\
\hline $\;\;$ & 4 & 6 & $\;\;$ & $\;\;$ & $\;\;$ \\
\hline 6 & $\;\;$ & 2 & $\;\;$ & $\;\;$ & $\;\;$ \\
\hline 
\end{tabular}\\ 
6.3, size 16
 & 6.3, size 17
 & 6.4, size 10
 & 6.4, size 11
\\
\end{tabular} 
\end{center} 
\begin{center}
\begin{tabular}{c@{\hspace{1.0mm}}c@{\hspace{1.0mm}}c@{\hspace{1.0mm}}c@{\hspace{1.0mm}}}
\begin{tabular}{|c|c|c|c|c|c|}
\hline $\;\;$ & $\;\;$ & $\;\;$ & $\;\;$ & $\;\;$ & $\;\;$ \\
\hline $\;\;$ & 1 & $\;\;$ & 3 & $\;\;$ & 5 \\
\hline $\;\;$ & $\;\;$ & 1 & $\;\;$ & $\;\;$ & 2 \\
\hline $\;\;$ & $\;\;$ & $\;\;$ & 1 & $\;\;$ & 3 \\
\hline $\;\;$ & 4 & 6 & $\;\;$ & $\;\;$ & $\;\;$ \\
\hline 6 & 3 & $\;\;$ & 5 & $\;\;$ & $\;\;$ \\
\hline 
\end{tabular} &
\begin{tabular}{|c|c|c|c|c|c|}
\hline $\;\;$ & $\;\;$ & $\;\;$ & $\;\;$ & $\;\;$ & $\;\;$ \\
\hline $\;\;$ & 1 & $\;\;$ & 3 & $\;\;$ & 5 \\
\hline $\;\;$ & $\;\;$ & 1 & $\;\;$ & $\;\;$ & 2 \\
\hline $\;\;$ & $\;\;$ & $\;\;$ & 1 & $\;\;$ & 3 \\
\hline $\;\;$ & $\;\;$ & 6 & 2 & 3 & $\;\;$ \\
\hline 6 & 3 & $\;\;$ & $\;\;$ & 1 & $\;\;$ \\
\hline 
\end{tabular} &
\begin{tabular}{|c|c|c|c|c|c|}
\hline $\;\;$ & $\;\;$ & $\;\;$ & $\;\;$ & $\;\;$ & $\;\;$ \\
\hline $\;\;$ & 1 & $\;\;$ & 3 & $\;\;$ & 5 \\
\hline $\;\;$ & $\;\;$ & 1 & $\;\;$ & $\;\;$ & 2 \\
\hline $\;\;$ & $\;\;$ & $\;\;$ & 1 & $\;\;$ & 3 \\
\hline $\;\;$ & $\;\;$ & 6 & 2 & 3 & $\;\;$ \\
\hline $\;\;$ & 3 & $\;\;$ & 5 & 1 & 4 \\
\hline 
\end{tabular} &
\begin{tabular}{|c|c|c|c|c|c|}
\hline $\;\;$ & $\;\;$ & $\;\;$ & $\;\;$ & $\;\;$ & $\;\;$ \\
\hline $\;\;$ & 1 & $\;\;$ & 3 & $\;\;$ & 5 \\
\hline $\;\;$ & $\;\;$ & 1 & $\;\;$ & $\;\;$ & 2 \\
\hline $\;\;$ & $\;\;$ & $\;\;$ & 1 & $\;\;$ & 3 \\
\hline $\;\;$ & 4 & 6 & 2 & 3 & $\;\;$ \\
\hline $\;\;$ & 3 & 2 & $\;\;$ & 1 & 4 \\
\hline 
\end{tabular}\\ 
6.4, size 12
 & 6.4, size 13
 & 6.4, size 14
 & 6.4, size 15
\\
\end{tabular} 
\end{center} 
\begin{center}
\begin{tabular}{c@{\hspace{1.0mm}}c@{\hspace{1.0mm}}c@{\hspace{1.0mm}}c@{\hspace{1.0mm}}}
\begin{tabular}{|c|c|c|c|c|c|}
\hline $\;\;$ & $\;\;$ & $\;\;$ & $\;\;$ & $\;\;$ & $\;\;$ \\
\hline $\;\;$ & 1 & $\;\;$ & 3 & $\;\;$ & 5 \\
\hline $\;\;$ & $\;\;$ & 1 & $\;\;$ & $\;\;$ & 2 \\
\hline $\;\;$ & $\;\;$ & $\;\;$ & 1 & $\;\;$ & 3 \\
\hline $\;\;$ & 4 & $\;\;$ & 2 & 3 & 1 \\
\hline $\;\;$ & 3 & 2 & 5 & 1 & 4 \\
\hline 
\end{tabular} &
\begin{tabular}{|c|c|c|c|c|c|}
\hline $\;\;$ & $\;\;$ & $\;\;$ & $\;\;$ & $\;\;$ & $\;\;$ \\
\hline $\;\;$ & 1 & $\;\;$ & 3 & $\;\;$ & 5 \\
\hline 3 & $\;\;$ & 1 & $\;\;$ & $\;\;$ & 2 \\
\hline 4 & $\;\;$ & 5 & 1 & $\;\;$ & 3 \\
\hline 5 & 4 & $\;\;$ & 2 & $\;\;$ & $\;\;$ \\
\hline $\;\;$ & 3 & 2 & 5 & $\;\;$ & 4 \\
\hline 
\end{tabular} &
\begin{tabular}{|c|c|c|c|c|c|}
\hline $\;\;$ & $\;\;$ & $\;\;$ & $\;\;$ & $\;\;$ & 6 \\
\hline $\;\;$ & 1 & $\;\;$ & 3 & $\;\;$ & 5 \\
\hline $\;\;$ & $\;\;$ & 5 & $\;\;$ & $\;\;$ & $\;\;$ \\
\hline $\;\;$ & 5 & $\;\;$ & $\;\;$ & $\;\;$ & 1 \\
\hline $\;\;$ & $\;\;$ & 2 & $\;\;$ & 4 & $\;\;$ \\
\hline $\;\;$ & $\;\;$ & $\;\;$ & $\;\;$ & 2 & $\;\;$ \\
\hline 
\end{tabular} &
\begin{tabular}{|c|c|c|c|c|c|}
\hline $\;\;$ & $\;\;$ & $\;\;$ & $\;\;$ & $\;\;$ & $\;\;$ \\
\hline $\;\;$ & 1 & $\;\;$ & 3 & $\;\;$ & 5 \\
\hline $\;\;$ & $\;\;$ & $\;\;$ & $\;\;$ & 1 & 2 \\
\hline $\;\;$ & 5 & $\;\;$ & $\;\;$ & $\;\;$ & $\;\;$ \\
\hline $\;\;$ & $\;\;$ & 2 & $\;\;$ & 4 & $\;\;$ \\
\hline $\;\;$ & $\;\;$ & 1 & $\;\;$ & 2 & 4 \\
\hline 
\end{tabular}\\ 
6.4, size 16
 & 6.4, size 17
 & 6.5, size 10
 & 6.5, size 11
\\
\end{tabular} 
\end{center} 
\begin{center}
\begin{tabular}{c@{\hspace{1.0mm}}c@{\hspace{1.0mm}}c@{\hspace{1.0mm}}c@{\hspace{1.0mm}}}
\begin{tabular}{|c|c|c|c|c|c|}
\hline $\;\;$ & $\;\;$ & $\;\;$ & $\;\;$ & $\;\;$ & $\;\;$ \\
\hline $\;\;$ & 1 & $\;\;$ & 3 & $\;\;$ & 5 \\
\hline $\;\;$ & $\;\;$ & $\;\;$ & $\;\;$ & 1 & 2 \\
\hline $\;\;$ & $\;\;$ & $\;\;$ & 2 & 3 & 1 \\
\hline 5 & 6 & $\;\;$ & $\;\;$ & $\;\;$ & $\;\;$ \\
\hline 6 & $\;\;$ & $\;\;$ & $\;\;$ & 2 & $\;\;$ \\
\hline 
\end{tabular} &
\begin{tabular}{|c|c|c|c|c|c|}
\hline $\;\;$ & $\;\;$ & $\;\;$ & $\;\;$ & $\;\;$ & $\;\;$ \\
\hline $\;\;$ & 1 & $\;\;$ & 3 & $\;\;$ & 5 \\
\hline $\;\;$ & $\;\;$ & $\;\;$ & $\;\;$ & 1 & 2 \\
\hline $\;\;$ & $\;\;$ & $\;\;$ & 2 & 3 & 1 \\
\hline $\;\;$ & $\;\;$ & 2 & 1 & 4 & 3 \\
\hline 6 & $\;\;$ & $\;\;$ & $\;\;$ & $\;\;$ & $\;\;$ \\
\hline 
\end{tabular} &
\begin{tabular}{|c|c|c|c|c|c|}
\hline $\;\;$ & $\;\;$ & $\;\;$ & $\;\;$ & $\;\;$ & $\;\;$ \\
\hline $\;\;$ & 1 & $\;\;$ & 3 & $\;\;$ & 5 \\
\hline $\;\;$ & $\;\;$ & $\;\;$ & $\;\;$ & 1 & 2 \\
\hline $\;\;$ & $\;\;$ & $\;\;$ & 2 & 3 & 1 \\
\hline 5 & $\;\;$ & $\;\;$ & $\;\;$ & 4 & 3 \\
\hline 6 & 3 & $\;\;$ & $\;\;$ & 2 & $\;\;$ \\
\hline 
\end{tabular} &
\begin{tabular}{|c|c|c|c|c|c|}
\hline $\;\;$ & $\;\;$ & $\;\;$ & $\;\;$ & $\;\;$ & $\;\;$ \\
\hline $\;\;$ & 1 & $\;\;$ & 3 & $\;\;$ & 5 \\
\hline $\;\;$ & $\;\;$ & $\;\;$ & $\;\;$ & 1 & 2 \\
\hline $\;\;$ & $\;\;$ & $\;\;$ & 2 & 3 & 1 \\
\hline $\;\;$ & 6 & 2 & $\;\;$ & 4 & $\;\;$ \\
\hline $\;\;$ & $\;\;$ & 1 & 5 & 2 & 4 \\
\hline 
\end{tabular}\\ 
6.5, size 12
 & 6.5, size 13
 & 6.5, size 14
 & 6.5, size 15
\\
\end{tabular} 
\end{center} 
\begin{center}
\begin{tabular}{c@{\hspace{1.0mm}}c@{\hspace{1.0mm}}c@{\hspace{1.0mm}}c@{\hspace{1.0mm}}}
\begin{tabular}{|c|c|c|c|c|c|}
\hline $\;\;$ & $\;\;$ & $\;\;$ & $\;\;$ & $\;\;$ & $\;\;$ \\
\hline $\;\;$ & 1 & $\;\;$ & 3 & $\;\;$ & 5 \\
\hline $\;\;$ & $\;\;$ & $\;\;$ & $\;\;$ & 1 & 2 \\
\hline $\;\;$ & $\;\;$ & $\;\;$ & 2 & 3 & 1 \\
\hline $\;\;$ & 6 & 2 & 1 & 4 & $\;\;$ \\
\hline $\;\;$ & 3 & 1 & $\;\;$ & 2 & 4 \\
\hline 
\end{tabular} &
\begin{tabular}{|c|c|c|c|c|c|}
\hline $\;\;$ & $\;\;$ & $\;\;$ & $\;\;$ & $\;\;$ & $\;\;$ \\
\hline $\;\;$ & 1 & $\;\;$ & 3 & $\;\;$ & 5 \\
\hline $\;\;$ & $\;\;$ & $\;\;$ & $\;\;$ & 1 & 2 \\
\hline $\;\;$ & $\;\;$ & $\;\;$ & 2 & 3 & 1 \\
\hline $\;\;$ & $\;\;$ & 2 & 1 & 4 & 3 \\
\hline $\;\;$ & 3 & 1 & 5 & 2 & 4 \\
\hline 
\end{tabular} &
\begin{tabular}{|c|c|c|c|c|c|}
\hline $\;\;$ & $\;\;$ & $\;\;$ & $\;\;$ & $\;\;$ & 6 \\
\hline $\;\;$ & 1 & $\;\;$ & 3 & $\;\;$ & $\;\;$ \\
\hline $\;\;$ & $\;\;$ & $\;\;$ & $\;\;$ & 1 & $\;\;$ \\
\hline 4 & 5 & $\;\;$ & $\;\;$ & $\;\;$ & $\;\;$ \\
\hline 5 & 6 & $\;\;$ & $\;\;$ & $\;\;$ & 4 \\
\hline $\;\;$ & $\;\;$ & 2 & $\;\;$ & 4 & $\;\;$ \\
\hline 
\end{tabular} &
\begin{tabular}{|c|c|c|c|c|c|}
\hline $\;\;$ & $\;\;$ & $\;\;$ & $\;\;$ & $\;\;$ & 6 \\
\hline $\;\;$ & 1 & $\;\;$ & 3 & $\;\;$ & $\;\;$ \\
\hline $\;\;$ & $\;\;$ & $\;\;$ & $\;\;$ & 1 & $\;\;$ \\
\hline $\;\;$ & $\;\;$ & $\;\;$ & 1 & 2 & 3 \\
\hline 5 & $\;\;$ & $\;\;$ & 2 & $\;\;$ & $\;\;$ \\
\hline 6 & $\;\;$ & 2 & $\;\;$ & 4 & $\;\;$ \\
\hline 
\end{tabular}\\ 
6.5, size 16
 & 6.5, size 17
 & 6.6, size 11
 & 6.6, size 12
\\
\end{tabular} 
\end{center} 
\begin{center}
\begin{tabular}{c@{\hspace{1.0mm}}c@{\hspace{1.0mm}}c@{\hspace{1.0mm}}c@{\hspace{1.0mm}}}
\begin{tabular}{|c|c|c|c|c|c|}
\hline $\;\;$ & $\;\;$ & $\;\;$ & $\;\;$ & $\;\;$ & $\;\;$ \\
\hline $\;\;$ & 1 & $\;\;$ & 3 & $\;\;$ & 5 \\
\hline $\;\;$ & $\;\;$ & $\;\;$ & $\;\;$ & 1 & 2 \\
\hline $\;\;$ & $\;\;$ & $\;\;$ & 1 & 2 & 3 \\
\hline 5 & 6 & $\;\;$ & $\;\;$ & $\;\;$ & $\;\;$ \\
\hline 6 & $\;\;$ & 2 & $\;\;$ & 4 & $\;\;$ \\
\hline 
\end{tabular} &
\begin{tabular}{|c|c|c|c|c|c|}
\hline $\;\;$ & $\;\;$ & $\;\;$ & $\;\;$ & $\;\;$ & $\;\;$ \\
\hline $\;\;$ & 1 & $\;\;$ & 3 & $\;\;$ & 5 \\
\hline $\;\;$ & $\;\;$ & $\;\;$ & $\;\;$ & 1 & 2 \\
\hline $\;\;$ & $\;\;$ & $\;\;$ & 1 & 2 & 3 \\
\hline $\;\;$ & 6 & 1 & $\;\;$ & 3 & $\;\;$ \\
\hline 6 & $\;\;$ & 2 & $\;\;$ & 4 & $\;\;$ \\
\hline 
\end{tabular} &
\begin{tabular}{|c|c|c|c|c|c|}
\hline $\;\;$ & $\;\;$ & $\;\;$ & $\;\;$ & $\;\;$ & $\;\;$ \\
\hline $\;\;$ & 1 & $\;\;$ & 3 & $\;\;$ & 5 \\
\hline $\;\;$ & $\;\;$ & $\;\;$ & $\;\;$ & 1 & 2 \\
\hline $\;\;$ & $\;\;$ & $\;\;$ & 1 & 2 & 3 \\
\hline $\;\;$ & $\;\;$ & 1 & 2 & 3 & 4 \\
\hline 6 & $\;\;$ & 2 & $\;\;$ & 4 & $\;\;$ \\
\hline 
\end{tabular} &
\begin{tabular}{|c|c|c|c|c|c|}
\hline $\;\;$ & $\;\;$ & $\;\;$ & $\;\;$ & $\;\;$ & $\;\;$ \\
\hline $\;\;$ & 1 & $\;\;$ & 3 & $\;\;$ & 5 \\
\hline $\;\;$ & $\;\;$ & $\;\;$ & $\;\;$ & 1 & 2 \\
\hline $\;\;$ & $\;\;$ & 6 & 1 & 2 & $\;\;$ \\
\hline $\;\;$ & $\;\;$ & 1 & 2 & 3 & 4 \\
\hline $\;\;$ & 3 & $\;\;$ & 5 & 4 & 1 \\
\hline 
\end{tabular}\\ 
6.6, size 13
 & 6.6, size 14
 & 6.6, size 15
 & 6.6, size 16
\\
\end{tabular} 
\end{center} 
\begin{center}
\begin{tabular}{c@{\hspace{1.0mm}}c@{\hspace{1.0mm}}c@{\hspace{1.0mm}}c@{\hspace{1.0mm}}}
\begin{tabular}{|c|c|c|c|c|c|}
\hline $\;\;$ & $\;\;$ & $\;\;$ & $\;\;$ & $\;\;$ & $\;\;$ \\
\hline $\;\;$ & 1 & $\;\;$ & 3 & $\;\;$ & 5 \\
\hline $\;\;$ & $\;\;$ & $\;\;$ & $\;\;$ & 1 & 2 \\
\hline $\;\;$ & $\;\;$ & $\;\;$ & 1 & 2 & 3 \\
\hline $\;\;$ & $\;\;$ & 1 & 2 & 3 & 4 \\
\hline $\;\;$ & 3 & 2 & 5 & 4 & 1 \\
\hline 
\end{tabular} &
\begin{tabular}{|c|c|c|c|c|c|}
\hline $\;\;$ & $\;\;$ & $\;\;$ & $\;\;$ & $\;\;$ & 6 \\
\hline $\;\;$ & $\;\;$ & 1 & $\;\;$ & $\;\;$ & $\;\;$ \\
\hline $\;\;$ & 1 & 2 & 5 & $\;\;$ & $\;\;$ \\
\hline 4 & 5 & $\;\;$ & $\;\;$ & $\;\;$ & $\;\;$ \\
\hline 5 & $\;\;$ & $\;\;$ & 3 & $\;\;$ & $\;\;$ \\
\hline $\;\;$ & $\;\;$ & $\;\;$ & 2 & 3 & 1 \\
\hline 
\end{tabular} &
\begin{tabular}{|c|c|c|c|c|c|}
\hline $\;\;$ & $\;\;$ & $\;\;$ & $\;\;$ & $\;\;$ & $\;\;$ \\
\hline $\;\;$ & $\;\;$ & 1 & $\;\;$ & $\;\;$ & 5 \\
\hline 3 & $\;\;$ & $\;\;$ & $\;\;$ & 6 & $\;\;$ \\
\hline $\;\;$ & $\;\;$ & $\;\;$ & 1 & 2 & 3 \\
\hline $\;\;$ & 6 & 4 & $\;\;$ & 1 & $\;\;$ \\
\hline 6 & 4 & 5 & $\;\;$ & $\;\;$ & $\;\;$ \\
\hline 
\end{tabular} &
\begin{tabular}{|c|c|c|c|c|c|}
\hline $\;\;$ & $\;\;$ & $\;\;$ & $\;\;$ & $\;\;$ & $\;\;$ \\
\hline $\;\;$ & $\;\;$ & 1 & $\;\;$ & $\;\;$ & 5 \\
\hline $\;\;$ & 1 & 2 & $\;\;$ & 6 & $\;\;$ \\
\hline $\;\;$ & $\;\;$ & $\;\;$ & 1 & 2 & 3 \\
\hline $\;\;$ & 6 & $\;\;$ & 3 & 1 & $\;\;$ \\
\hline 6 & 4 & $\;\;$ & $\;\;$ & 3 & $\;\;$ \\
\hline 
\end{tabular}\\ 
6.6, size 17
 & 6.7, size 12
 & 6.7, size 13
 & 6.7, size 14
\\
\end{tabular} 
\end{center} 
\begin{center}
\begin{tabular}{c@{\hspace{1.0mm}}c@{\hspace{1.0mm}}c@{\hspace{1.0mm}}c@{\hspace{1.0mm}}}
\begin{tabular}{|c|c|c|c|c|c|}
\hline $\;\;$ & $\;\;$ & $\;\;$ & $\;\;$ & $\;\;$ & $\;\;$ \\
\hline $\;\;$ & $\;\;$ & 1 & $\;\;$ & $\;\;$ & 5 \\
\hline $\;\;$ & 1 & 2 & $\;\;$ & 6 & $\;\;$ \\
\hline $\;\;$ & $\;\;$ & $\;\;$ & 1 & 2 & 3 \\
\hline $\;\;$ & $\;\;$ & 4 & 3 & 1 & 2 \\
\hline 6 & $\;\;$ & $\;\;$ & 2 & 3 & $\;\;$ \\
\hline 
\end{tabular} &
\begin{tabular}{|c|c|c|c|c|c|}
\hline $\;\;$ & $\;\;$ & $\;\;$ & $\;\;$ & $\;\;$ & $\;\;$ \\
\hline $\;\;$ & $\;\;$ & 1 & $\;\;$ & $\;\;$ & 5 \\
\hline $\;\;$ & 1 & 2 & $\;\;$ & 6 & $\;\;$ \\
\hline $\;\;$ & $\;\;$ & $\;\;$ & 1 & 2 & 3 \\
\hline $\;\;$ & $\;\;$ & 4 & 3 & 1 & 2 \\
\hline $\;\;$ & 4 & 5 & 2 & $\;\;$ & 1 \\
\hline 
\end{tabular} &
\begin{tabular}{|c|c|c|c|c|c|}
\hline $\;\;$ & $\;\;$ & $\;\;$ & $\;\;$ & $\;\;$ & $\;\;$ \\
\hline $\;\;$ & $\;\;$ & 1 & $\;\;$ & $\;\;$ & 5 \\
\hline $\;\;$ & 1 & 2 & $\;\;$ & 6 & $\;\;$ \\
\hline $\;\;$ & 5 & 6 & 1 & 2 & 3 \\
\hline $\;\;$ & $\;\;$ & $\;\;$ & 3 & 1 & 2 \\
\hline $\;\;$ & $\;\;$ & 5 & 2 & 3 & 1 \\
\hline 
\end{tabular} &
\begin{tabular}{|c|c|c|c|c|c|}
\hline $\;\;$ & $\;\;$ & $\;\;$ & $\;\;$ & $\;\;$ & $\;\;$ \\
\hline $\;\;$ & $\;\;$ & 1 & $\;\;$ & $\;\;$ & 5 \\
\hline $\;\;$ & 1 & 2 & 5 & $\;\;$ & 4 \\
\hline $\;\;$ & $\;\;$ & $\;\;$ & 1 & 2 & 3 \\
\hline $\;\;$ & $\;\;$ & 4 & 3 & 1 & 2 \\
\hline $\;\;$ & 4 & 5 & 2 & 3 & 1 \\
\hline 
\end{tabular}\\ 
6.7, size 15
 & 6.7, size 16
 & 6.7, size 17
 & 6.7, size 18
\\
\end{tabular} 
\end{center} 
\begin{center}
\begin{tabular}{c@{\hspace{1.0mm}}c@{\hspace{1.0mm}}c@{\hspace{1.0mm}}c@{\hspace{1.0mm}}}
\begin{tabular}{|c|c|c|c|c|c|}
\hline $\;\;$ & $\;\;$ & $\;\;$ & $\;\;$ & $\;\;$ & $\;\;$ \\
\hline $\;\;$ & 1 & $\;\;$ & 3 & 6 & $\;\;$ \\
\hline $\;\;$ & 5 & 1 & $\;\;$ & $\;\;$ & 4 \\
\hline 4 & $\;\;$ & $\;\;$ & $\;\;$ & $\;\;$ & 3 \\
\hline $\;\;$ & $\;\;$ & 6 & $\;\;$ & $\;\;$ & $\;\;$ \\
\hline $\;\;$ & $\;\;$ & $\;\;$ & $\;\;$ & 3 & 2 \\
\hline 
\end{tabular} &
\begin{tabular}{|c|c|c|c|c|c|}
\hline $\;\;$ & $\;\;$ & $\;\;$ & $\;\;$ & $\;\;$ & $\;\;$ \\
\hline $\;\;$ & 1 & $\;\;$ & 3 & $\;\;$ & 5 \\
\hline $\;\;$ & $\;\;$ & 1 & $\;\;$ & 2 & 4 \\
\hline 4 & 6 & $\;\;$ & $\;\;$ & $\;\;$ & $\;\;$ \\
\hline $\;\;$ & 3 & 6 & $\;\;$ & $\;\;$ & $\;\;$ \\
\hline $\;\;$ & $\;\;$ & $\;\;$ & $\;\;$ & 3 & 2 \\
\hline 
\end{tabular} &
\begin{tabular}{|c|c|c|c|c|c|}
\hline $\;\;$ & $\;\;$ & $\;\;$ & $\;\;$ & $\;\;$ & $\;\;$ \\
\hline $\;\;$ & 1 & $\;\;$ & 3 & $\;\;$ & 5 \\
\hline $\;\;$ & $\;\;$ & 1 & $\;\;$ & 2 & 4 \\
\hline $\;\;$ & $\;\;$ & 2 & 5 & $\;\;$ & $\;\;$ \\
\hline $\;\;$ & 3 & $\;\;$ & $\;\;$ & 4 & $\;\;$ \\
\hline $\;\;$ & 4 & 5 & $\;\;$ & $\;\;$ & 2 \\
\hline 
\end{tabular} &
\begin{tabular}{|c|c|c|c|c|c|}
\hline $\;\;$ & $\;\;$ & $\;\;$ & $\;\;$ & $\;\;$ & $\;\;$ \\
\hline $\;\;$ & 1 & $\;\;$ & 3 & $\;\;$ & 5 \\
\hline $\;\;$ & $\;\;$ & 1 & $\;\;$ & 2 & 4 \\
\hline $\;\;$ & $\;\;$ & 2 & 5 & $\;\;$ & $\;\;$ \\
\hline $\;\;$ & 3 & 6 & $\;\;$ & 4 & $\;\;$ \\
\hline $\;\;$ & 4 & $\;\;$ & $\;\;$ & 3 & 2 \\
\hline 
\end{tabular}\\ 
6.8, size 11
 & 6.8, size 12
 & 6.8, size 13
 & 6.8, size 14
\\
\end{tabular} 
\end{center} 
\begin{center}
\begin{tabular}{c@{\hspace{1.0mm}}c@{\hspace{1.0mm}}c@{\hspace{1.0mm}}c@{\hspace{1.0mm}}}
\begin{tabular}{|c|c|c|c|c|c|}
\hline $\;\;$ & $\;\;$ & $\;\;$ & $\;\;$ & $\;\;$ & $\;\;$ \\
\hline $\;\;$ & 1 & $\;\;$ & 3 & $\;\;$ & 5 \\
\hline $\;\;$ & $\;\;$ & 1 & $\;\;$ & 2 & 4 \\
\hline $\;\;$ & $\;\;$ & 2 & 5 & $\;\;$ & $\;\;$ \\
\hline 5 & $\;\;$ & 6 & $\;\;$ & 4 & 1 \\
\hline 6 & $\;\;$ & 5 & $\;\;$ & $\;\;$ & 2 \\
\hline 
\end{tabular} &
\begin{tabular}{|c|c|c|c|c|c|}
\hline $\;\;$ & $\;\;$ & $\;\;$ & $\;\;$ & $\;\;$ & $\;\;$ \\
\hline $\;\;$ & 1 & $\;\;$ & 3 & $\;\;$ & 5 \\
\hline $\;\;$ & $\;\;$ & 1 & 6 & 2 & $\;\;$ \\
\hline $\;\;$ & 6 & 2 & $\;\;$ & 1 & 3 \\
\hline $\;\;$ & 3 & 6 & $\;\;$ & 4 & 1 \\
\hline $\;\;$ & $\;\;$ & $\;\;$ & 1 & 3 & $\;\;$ \\
\hline 
\end{tabular} &
\begin{tabular}{|c|c|c|c|c|c|}
\hline $\;\;$ & $\;\;$ & $\;\;$ & $\;\;$ & $\;\;$ & $\;\;$ \\
\hline $\;\;$ & 1 & $\;\;$ & 3 & $\;\;$ & 5 \\
\hline 3 & 5 & $\;\;$ & $\;\;$ & $\;\;$ & 4 \\
\hline 4 & $\;\;$ & $\;\;$ & 5 & 1 & 3 \\
\hline 5 & 3 & $\;\;$ & $\;\;$ & 4 & 1 \\
\hline $\;\;$ & 4 & $\;\;$ & $\;\;$ & 3 & 2 \\
\hline 
\end{tabular} &
\begin{tabular}{|c|c|c|c|c|c|}
\hline $\;\;$ & $\;\;$ & $\;\;$ & $\;\;$ & $\;\;$ & 6 \\
\hline $\;\;$ & 1 & 4 & 3 & $\;\;$ & $\;\;$ \\
\hline $\;\;$ & 5 & 1 & $\;\;$ & $\;\;$ & $\;\;$ \\
\hline $\;\;$ & $\;\;$ & $\;\;$ & $\;\;$ & $\;\;$ & $\;\;$ \\
\hline $\;\;$ & 4 & $\;\;$ & $\;\;$ & 1 & $\;\;$ \\
\hline 6 & $\;\;$ & $\;\;$ & $\;\;$ & $\;\;$ & 2 \\
\hline 
\end{tabular}\\ 
6.8, size 15
 & 6.8, size 16
 & 6.8, size 17
 & 6.9, size 10
\\
\end{tabular} 
\end{center} 
\begin{center}
\begin{tabular}{c@{\hspace{1.0mm}}c@{\hspace{1.0mm}}c@{\hspace{1.0mm}}c@{\hspace{1.0mm}}}
\begin{tabular}{|c|c|c|c|c|c|}
\hline $\;\;$ & $\;\;$ & $\;\;$ & $\;\;$ & $\;\;$ & $\;\;$ \\
\hline $\;\;$ & 1 & $\;\;$ & 3 & $\;\;$ & 5 \\
\hline $\;\;$ & $\;\;$ & 1 & 6 & 2 & $\;\;$ \\
\hline $\;\;$ & $\;\;$ & 2 & $\;\;$ & $\;\;$ & $\;\;$ \\
\hline 5 & 4 & $\;\;$ & $\;\;$ & $\;\;$ & 3 \\
\hline $\;\;$ & 3 & $\;\;$ & $\;\;$ & $\;\;$ & $\;\;$ \\
\hline 
\end{tabular} &
\begin{tabular}{|c|c|c|c|c|c|}
\hline $\;\;$ & $\;\;$ & $\;\;$ & $\;\;$ & $\;\;$ & $\;\;$ \\
\hline $\;\;$ & 1 & $\;\;$ & 3 & $\;\;$ & 5 \\
\hline $\;\;$ & $\;\;$ & 1 & $\;\;$ & 2 & 4 \\
\hline $\;\;$ & $\;\;$ & 2 & 5 & $\;\;$ & $\;\;$ \\
\hline $\;\;$ & $\;\;$ & 6 & 2 & 1 & $\;\;$ \\
\hline $\;\;$ & 3 & $\;\;$ & $\;\;$ & $\;\;$ & $\;\;$ \\
\hline 
\end{tabular} &
\begin{tabular}{|c|c|c|c|c|c|}
\hline $\;\;$ & $\;\;$ & $\;\;$ & $\;\;$ & $\;\;$ & $\;\;$ \\
\hline $\;\;$ & 1 & $\;\;$ & 3 & $\;\;$ & 5 \\
\hline $\;\;$ & $\;\;$ & 1 & $\;\;$ & 2 & 4 \\
\hline $\;\;$ & $\;\;$ & 2 & 5 & $\;\;$ & $\;\;$ \\
\hline $\;\;$ & $\;\;$ & 6 & 2 & 1 & $\;\;$ \\
\hline 6 & $\;\;$ & $\;\;$ & $\;\;$ & 4 & $\;\;$ \\
\hline 
\end{tabular} &
\begin{tabular}{|c|c|c|c|c|c|}
\hline $\;\;$ & $\;\;$ & $\;\;$ & $\;\;$ & $\;\;$ & $\;\;$ \\
\hline $\;\;$ & 1 & $\;\;$ & 3 & $\;\;$ & 5 \\
\hline $\;\;$ & $\;\;$ & 1 & $\;\;$ & 2 & 4 \\
\hline $\;\;$ & $\;\;$ & 2 & 5 & $\;\;$ & $\;\;$ \\
\hline $\;\;$ & $\;\;$ & $\;\;$ & 2 & 1 & 3 \\
\hline $\;\;$ & 3 & 5 & $\;\;$ & 4 & $\;\;$ \\
\hline 
\end{tabular}\\ 
6.9, size 11
 & 6.9, size 12
 & 6.9, size 13
 & 6.9, size 14
\\
\end{tabular} 
\end{center} 
\begin{center}
\begin{tabular}{c@{\hspace{1.0mm}}c@{\hspace{1.0mm}}c@{\hspace{1.0mm}}c@{\hspace{1.0mm}}}
\begin{tabular}{|c|c|c|c|c|c|}
\hline $\;\;$ & $\;\;$ & $\;\;$ & $\;\;$ & $\;\;$ & $\;\;$ \\
\hline $\;\;$ & 1 & $\;\;$ & 3 & $\;\;$ & 5 \\
\hline $\;\;$ & $\;\;$ & 1 & $\;\;$ & 2 & 4 \\
\hline 4 & 6 & $\;\;$ & $\;\;$ & $\;\;$ & $\;\;$ \\
\hline 5 & 4 & 6 & $\;\;$ & $\;\;$ & 3 \\
\hline 6 & 3 & 5 & $\;\;$ & $\;\;$ & $\;\;$ \\
\hline 
\end{tabular} &
\begin{tabular}{|c|c|c|c|c|c|}
\hline $\;\;$ & $\;\;$ & $\;\;$ & $\;\;$ & $\;\;$ & $\;\;$ \\
\hline $\;\;$ & 1 & $\;\;$ & 3 & $\;\;$ & 5 \\
\hline 3 & $\;\;$ & $\;\;$ & $\;\;$ & 2 & 4 \\
\hline 4 & $\;\;$ & $\;\;$ & 5 & 3 & 1 \\
\hline 5 & 4 & $\;\;$ & 2 & 1 & 3 \\
\hline $\;\;$ & $\;\;$ & $\;\;$ & $\;\;$ & $\;\;$ & 2 \\
\hline 
\end{tabular} &
\begin{tabular}{|c|c|c|c|c|c|}
\hline $\;\;$ & 2 & 3 & 4 & 5 & 6 \\
\hline $\;\;$ & $\;\;$ & 4 & 3 & $\;\;$ & 5 \\
\hline $\;\;$ & 5 & $\;\;$ & $\;\;$ & $\;\;$ & 4 \\
\hline $\;\;$ & $\;\;$ & $\;\;$ & $\;\;$ & $\;\;$ & $\;\;$ \\
\hline $\;\;$ & 4 & 6 & $\;\;$ & $\;\;$ & 3 \\
\hline $\;\;$ & 3 & 5 & $\;\;$ & 4 & 2 \\
\hline 
\end{tabular} &
\begin{tabular}{|c|c|c|c|c|c|}
\hline 1 & $\;\;$ & $\;\;$ & $\;\;$ & 5 & 6 \\
\hline $\;\;$ & $\;\;$ & $\;\;$ & 3 & $\;\;$ & $\;\;$ \\
\hline $\;\;$ & $\;\;$ & 1 & $\;\;$ & $\;\;$ & $\;\;$ \\
\hline 4 & $\;\;$ & $\;\;$ & $\;\;$ & $\;\;$ & $\;\;$ \\
\hline $\;\;$ & 3 & $\;\;$ & 2 & $\;\;$ & $\;\;$ \\
\hline 6 & $\;\;$ & $\;\;$ & $\;\;$ & $\;\;$ & 1 \\
\hline 
\end{tabular}\\ 
6.9, size 15
 & 6.9, size 16
 & 6.9, size 17
 & 6.10, size 10
\\
\end{tabular} 
\end{center} 
\begin{center}
\begin{tabular}{c@{\hspace{1.0mm}}c@{\hspace{1.0mm}}c@{\hspace{1.0mm}}c@{\hspace{1.0mm}}}
\begin{tabular}{|c|c|c|c|c|c|}
\hline $\;\;$ & $\;\;$ & $\;\;$ & $\;\;$ & $\;\;$ & 6 \\
\hline $\;\;$ & 1 & $\;\;$ & 3 & $\;\;$ & $\;\;$ \\
\hline $\;\;$ & $\;\;$ & 1 & $\;\;$ & $\;\;$ & $\;\;$ \\
\hline 4 & $\;\;$ & $\;\;$ & $\;\;$ & $\;\;$ & $\;\;$ \\
\hline $\;\;$ & 3 & $\;\;$ & 2 & 1 & $\;\;$ \\
\hline 6 & $\;\;$ & $\;\;$ & 5 & 3 & $\;\;$ \\
\hline 
\end{tabular} &
\begin{tabular}{|c|c|c|c|c|c|}
\hline $\;\;$ & $\;\;$ & $\;\;$ & $\;\;$ & $\;\;$ & $\;\;$ \\
\hline $\;\;$ & 1 & $\;\;$ & 3 & $\;\;$ & 5 \\
\hline $\;\;$ & $\;\;$ & 1 & $\;\;$ & $\;\;$ & 2 \\
\hline $\;\;$ & $\;\;$ & $\;\;$ & 1 & $\;\;$ & 3 \\
\hline 5 & $\;\;$ & 6 & $\;\;$ & $\;\;$ & $\;\;$ \\
\hline 6 & 4 & $\;\;$ & $\;\;$ & 3 & $\;\;$ \\
\hline 
\end{tabular} &
\begin{tabular}{|c|c|c|c|c|c|}
\hline $\;\;$ & $\;\;$ & $\;\;$ & $\;\;$ & $\;\;$ & $\;\;$ \\
\hline $\;\;$ & 1 & $\;\;$ & 3 & $\;\;$ & 5 \\
\hline $\;\;$ & $\;\;$ & 1 & $\;\;$ & $\;\;$ & 2 \\
\hline $\;\;$ & 6 & $\;\;$ & 1 & $\;\;$ & $\;\;$ \\
\hline $\;\;$ & 3 & 6 & $\;\;$ & 1 & $\;\;$ \\
\hline 6 & $\;\;$ & 2 & 5 & $\;\;$ & $\;\;$ \\
\hline 
\end{tabular} &
\begin{tabular}{|c|c|c|c|c|c|}
\hline $\;\;$ & $\;\;$ & $\;\;$ & $\;\;$ & $\;\;$ & $\;\;$ \\
\hline $\;\;$ & 1 & $\;\;$ & 3 & $\;\;$ & 5 \\
\hline $\;\;$ & $\;\;$ & 1 & $\;\;$ & $\;\;$ & 2 \\
\hline $\;\;$ & $\;\;$ & $\;\;$ & 1 & $\;\;$ & 3 \\
\hline $\;\;$ & 3 & $\;\;$ & 2 & 1 & 4 \\
\hline 6 & 4 & $\;\;$ & $\;\;$ & 3 & $\;\;$ \\
\hline 
\end{tabular}\\ 
6.10, size 11
 & 6.10, size 12
 & 6.10, size 13
 & 6.10, size 14
\\
\end{tabular} 
\end{center} 
\begin{center}
\begin{tabular}{c@{\hspace{1.0mm}}c@{\hspace{1.0mm}}c@{\hspace{1.0mm}}c@{\hspace{1.0mm}}}
\begin{tabular}{|c|c|c|c|c|c|}
\hline $\;\;$ & $\;\;$ & $\;\;$ & $\;\;$ & $\;\;$ & $\;\;$ \\
\hline $\;\;$ & 1 & $\;\;$ & 3 & $\;\;$ & 5 \\
\hline $\;\;$ & $\;\;$ & 1 & $\;\;$ & $\;\;$ & 2 \\
\hline $\;\;$ & $\;\;$ & $\;\;$ & 1 & $\;\;$ & 3 \\
\hline $\;\;$ & 3 & 6 & 2 & 1 & $\;\;$ \\
\hline $\;\;$ & 4 & 2 & $\;\;$ & 3 & 1 \\
\hline 
\end{tabular} &
\begin{tabular}{|c|c|c|c|c|c|}
\hline $\;\;$ & $\;\;$ & $\;\;$ & $\;\;$ & $\;\;$ & $\;\;$ \\
\hline $\;\;$ & 1 & $\;\;$ & 3 & $\;\;$ & 5 \\
\hline $\;\;$ & $\;\;$ & 1 & $\;\;$ & $\;\;$ & 2 \\
\hline $\;\;$ & $\;\;$ & $\;\;$ & 1 & $\;\;$ & 3 \\
\hline $\;\;$ & 3 & $\;\;$ & 2 & 1 & 4 \\
\hline $\;\;$ & 4 & 2 & 5 & 3 & 1 \\
\hline 
\end{tabular} &
\begin{tabular}{|c|c|c|c|c|c|}
\hline $\;\;$ & $\;\;$ & $\;\;$ & $\;\;$ & $\;\;$ & $\;\;$ \\
\hline $\;\;$ & 1 & $\;\;$ & 3 & 6 & $\;\;$ \\
\hline $\;\;$ & $\;\;$ & 1 & 6 & $\;\;$ & $\;\;$ \\
\hline $\;\;$ & 6 & 5 & 1 & 2 & 3 \\
\hline $\;\;$ & 3 & 6 & 2 & 1 & $\;\;$ \\
\hline $\;\;$ & $\;\;$ & $\;\;$ & 5 & 3 & 1 \\
\hline 
\end{tabular} &
\begin{tabular}{|c|c|c|c|c|c|}
\hline $\;\;$ & $\;\;$ & $\;\;$ & $\;\;$ & $\;\;$ & 6 \\
\hline 2 & 1 & $\;\;$ & $\;\;$ & $\;\;$ & $\;\;$ \\
\hline $\;\;$ & 4 & $\;\;$ & $\;\;$ & 1 & $\;\;$ \\
\hline $\;\;$ & $\;\;$ & $\;\;$ & $\;\;$ & 3 & $\;\;$ \\
\hline $\;\;$ & 3 & $\;\;$ & 1 & $\;\;$ & $\;\;$ \\
\hline 6 & $\;\;$ & $\;\;$ & $\;\;$ & $\;\;$ & 2 \\
\hline 
\end{tabular}\\ 
6.10, size 15
 & 6.10, size 16
 & 6.10, size 17
 & 6.11, size 10
\\
\end{tabular} 
\end{center} 
\begin{center}
\begin{tabular}{c@{\hspace{1.0mm}}c@{\hspace{1.0mm}}c@{\hspace{1.0mm}}c@{\hspace{1.0mm}}}
\begin{tabular}{|c|c|c|c|c|c|}
\hline $\;\;$ & $\;\;$ & $\;\;$ & $\;\;$ & $\;\;$ & $\;\;$ \\
\hline $\;\;$ & 1 & $\;\;$ & $\;\;$ & $\;\;$ & 3 \\
\hline $\;\;$ & $\;\;$ & $\;\;$ & $\;\;$ & 1 & 5 \\
\hline $\;\;$ & $\;\;$ & $\;\;$ & 2 & 3 & 1 \\
\hline 5 & $\;\;$ & 6 & $\;\;$ & $\;\;$ & $\;\;$ \\
\hline 6 & $\;\;$ & $\;\;$ & $\;\;$ & $\;\;$ & 2 \\
\hline 
\end{tabular} &
\begin{tabular}{|c|c|c|c|c|c|}
\hline $\;\;$ & $\;\;$ & $\;\;$ & $\;\;$ & $\;\;$ & $\;\;$ \\
\hline $\;\;$ & 1 & $\;\;$ & $\;\;$ & $\;\;$ & 3 \\
\hline $\;\;$ & $\;\;$ & $\;\;$ & $\;\;$ & 1 & 5 \\
\hline $\;\;$ & $\;\;$ & $\;\;$ & 2 & 3 & 1 \\
\hline $\;\;$ & $\;\;$ & 6 & 1 & 2 & $\;\;$ \\
\hline 6 & 5 & $\;\;$ & $\;\;$ & $\;\;$ & $\;\;$ \\
\hline 
\end{tabular} &
\begin{tabular}{|c|c|c|c|c|c|}
\hline $\;\;$ & $\;\;$ & $\;\;$ & $\;\;$ & $\;\;$ & $\;\;$ \\
\hline $\;\;$ & 1 & $\;\;$ & $\;\;$ & $\;\;$ & 3 \\
\hline $\;\;$ & $\;\;$ & $\;\;$ & $\;\;$ & 1 & 5 \\
\hline $\;\;$ & $\;\;$ & $\;\;$ & 2 & 3 & 1 \\
\hline 5 & $\;\;$ & $\;\;$ & $\;\;$ & 2 & 4 \\
\hline 6 & $\;\;$ & $\;\;$ & 3 & 4 & $\;\;$ \\
\hline 
\end{tabular} &
\begin{tabular}{|c|c|c|c|c|c|}
\hline $\;\;$ & $\;\;$ & $\;\;$ & $\;\;$ & $\;\;$ & $\;\;$ \\
\hline $\;\;$ & 1 & $\;\;$ & $\;\;$ & $\;\;$ & 3 \\
\hline $\;\;$ & $\;\;$ & $\;\;$ & $\;\;$ & 1 & 5 \\
\hline $\;\;$ & $\;\;$ & $\;\;$ & 2 & 3 & 1 \\
\hline 5 & $\;\;$ & $\;\;$ & $\;\;$ & 2 & 4 \\
\hline $\;\;$ & $\;\;$ & 1 & 3 & 4 & 2 \\
\hline 
\end{tabular}\\ 
6.11, size 11
 & 6.11, size 12
 & 6.11, size 13
 & 6.11, size 14
\\
\end{tabular} 
\end{center} 
\begin{center}
\begin{tabular}{c@{\hspace{1.0mm}}c@{\hspace{1.0mm}}c@{\hspace{1.0mm}}c@{\hspace{1.0mm}}}
\begin{tabular}{|c|c|c|c|c|c|}
\hline $\;\;$ & $\;\;$ & $\;\;$ & $\;\;$ & $\;\;$ & $\;\;$ \\
\hline $\;\;$ & 1 & $\;\;$ & $\;\;$ & $\;\;$ & 3 \\
\hline $\;\;$ & $\;\;$ & $\;\;$ & $\;\;$ & 1 & 5 \\
\hline $\;\;$ & $\;\;$ & $\;\;$ & 2 & 3 & 1 \\
\hline $\;\;$ & 3 & 6 & 1 & 2 & $\;\;$ \\
\hline $\;\;$ & $\;\;$ & 1 & 3 & 4 & 2 \\
\hline 
\end{tabular} &
\begin{tabular}{|c|c|c|c|c|c|}
\hline $\;\;$ & $\;\;$ & $\;\;$ & $\;\;$ & $\;\;$ & $\;\;$ \\
\hline $\;\;$ & 1 & $\;\;$ & $\;\;$ & $\;\;$ & 3 \\
\hline $\;\;$ & $\;\;$ & $\;\;$ & $\;\;$ & 1 & 5 \\
\hline $\;\;$ & $\;\;$ & $\;\;$ & 2 & 3 & 1 \\
\hline $\;\;$ & 3 & $\;\;$ & 1 & 2 & 4 \\
\hline $\;\;$ & 5 & 1 & 3 & 4 & 2 \\
\hline 
\end{tabular} &
\begin{tabular}{|c|c|c|c|c|c|}
\hline $\;\;$ & $\;\;$ & $\;\;$ & $\;\;$ & $\;\;$ & $\;\;$ \\
\hline $\;\;$ & 1 & $\;\;$ & 5 & 6 & $\;\;$ \\
\hline 3 & 4 & $\;\;$ & 6 & 1 & $\;\;$ \\
\hline $\;\;$ & 6 & $\;\;$ & $\;\;$ & 3 & $\;\;$ \\
\hline $\;\;$ & 3 & 6 & 1 & $\;\;$ & $\;\;$ \\
\hline 6 & 5 & 1 & 3 & 4 & $\;\;$ \\
\hline 
\end{tabular} &
\begin{tabular}{|c|c|c|c|c|c|}
\hline $\;\;$ & $\;\;$ & $\;\;$ & $\;\;$ & $\;\;$ & $\;\;$ \\
\hline $\;\;$ & $\;\;$ & 1 & $\;\;$ & $\;\;$ & 5 \\
\hline $\;\;$ & 1 & 2 & $\;\;$ & 6 & $\;\;$ \\
\hline $\;\;$ & $\;\;$ & $\;\;$ & $\;\;$ & 3 & $\;\;$ \\
\hline $\;\;$ & $\;\;$ & 4 & $\;\;$ & 2 & 3 \\
\hline 6 & $\;\;$ & $\;\;$ & 3 & $\;\;$ & $\;\;$ \\
\hline 
\end{tabular}\\ 
6.11, size 15
 & 6.11, size 16
 & 6.11, size 17
 & 6.12, size 11
\\
\end{tabular} 
\end{center} 
\begin{center}
\begin{tabular}{c@{\hspace{1.0mm}}c@{\hspace{1.0mm}}c@{\hspace{1.0mm}}c@{\hspace{1.0mm}}}
\begin{tabular}{|c|c|c|c|c|c|}
\hline $\;\;$ & $\;\;$ & $\;\;$ & $\;\;$ & $\;\;$ & $\;\;$ \\
\hline $\;\;$ & $\;\;$ & 1 & $\;\;$ & $\;\;$ & 5 \\
\hline $\;\;$ & 1 & 2 & $\;\;$ & 6 & $\;\;$ \\
\hline $\;\;$ & $\;\;$ & $\;\;$ & $\;\;$ & 3 & $\;\;$ \\
\hline $\;\;$ & $\;\;$ & 4 & $\;\;$ & 2 & 3 \\
\hline 6 & 4 & 5 & $\;\;$ & $\;\;$ & $\;\;$ \\
\hline 
\end{tabular} &
\begin{tabular}{|c|c|c|c|c|c|}
\hline $\;\;$ & $\;\;$ & $\;\;$ & $\;\;$ & $\;\;$ & $\;\;$ \\
\hline $\;\;$ & $\;\;$ & 1 & $\;\;$ & $\;\;$ & 5 \\
\hline $\;\;$ & 1 & 2 & $\;\;$ & 6 & $\;\;$ \\
\hline $\;\;$ & $\;\;$ & $\;\;$ & $\;\;$ & 3 & $\;\;$ \\
\hline $\;\;$ & $\;\;$ & 4 & 1 & 2 & 3 \\
\hline $\;\;$ & 4 & 5 & $\;\;$ & $\;\;$ & 2 \\
\hline 
\end{tabular} &
\begin{tabular}{|c|c|c|c|c|c|}
\hline $\;\;$ & $\;\;$ & $\;\;$ & $\;\;$ & $\;\;$ & $\;\;$ \\
\hline $\;\;$ & $\;\;$ & 1 & $\;\;$ & $\;\;$ & 5 \\
\hline $\;\;$ & 1 & 2 & $\;\;$ & 6 & $\;\;$ \\
\hline $\;\;$ & $\;\;$ & $\;\;$ & $\;\;$ & 3 & $\;\;$ \\
\hline $\;\;$ & 6 & 4 & $\;\;$ & 2 & 3 \\
\hline $\;\;$ & 4 & $\;\;$ & 3 & 1 & 2 \\
\hline 
\end{tabular} &
\begin{tabular}{|c|c|c|c|c|c|}
\hline $\;\;$ & $\;\;$ & $\;\;$ & $\;\;$ & $\;\;$ & $\;\;$ \\
\hline $\;\;$ & $\;\;$ & 1 & $\;\;$ & $\;\;$ & 5 \\
\hline $\;\;$ & 1 & 2 & $\;\;$ & 6 & $\;\;$ \\
\hline $\;\;$ & $\;\;$ & $\;\;$ & 2 & $\;\;$ & 1 \\
\hline $\;\;$ & $\;\;$ & 4 & 1 & 2 & 3 \\
\hline $\;\;$ & 4 & 5 & 3 & $\;\;$ & 2 \\
\hline 
\end{tabular}\\ 
6.12, size 12
 & 6.12, size 13
 & 6.12, size 14
 & 6.12, size 15
\\
\end{tabular} 
\end{center} 
\begin{center}
\begin{tabular}{c}
\begin{tabular}{|c|c|c|c|c|c|}
\hline $\;\;$ & $\;\;$ & $\;\;$ & $\;\;$ & $\;\;$ & $\;\;$ \\
\hline $\;\;$ & $\;\;$ & 1 & $\;\;$ & $\;\;$ & 5 \\
\hline $\;\;$ & 1 & 2 & $\;\;$ & 6 & 4 \\
\hline $\;\;$ & 5 & 6 & 2 & $\;\;$ & 1 \\
\hline 5 & 6 & 4 & 1 & $\;\;$ & $\;\;$ \\
\hline $\;\;$ & 4 & 5 & $\;\;$ & $\;\;$ & $\;\;$ \\
\hline 
\end{tabular}\\ 
6.12, size 16
\\
\end{tabular} 
\end{center} 


\section{Acknowledgements}
We would like to extend our sincere appreciation to Ian Wanless
for programs to determine isotopy and main classes in $O(n)$ time
and to Brendan McKay for proposing a simpler method of using
{\it nauty} ~\cite{nauty} to determine these classes.

\begin{thebibliography}{10}

\bibitem{AK2}
P.~Adams and A.~Khodkar.
\newblock Smallest critical sets for the latin squares of orders six and seven.
\newblock To appear in Journal of Combinatorial Mathematics and Combinatorial
  Computing.

\bibitem{MR2000k:05054}
P.~Adams and A.~Khodkar.
\newblock Smallest weak and smallest totally weak critical sets for the latin
  squares of order at most seven.
\newblock To appear in {\it Ars Combin}.

\bibitem{MR2000g:05034}
J.~A. Bate and G.~H.~J. van Rees.
\newblock The size of the smallest strong critical set in a {L}atin square.
\newblock {\em Ars Combin.}, {\bf 53} (1999), 73--83.

\bibitem{bandw}
D.~Bedford and D.~Whitehouse.
\newblock Products of uniquely completable partial latin squares.
\newblock To appear in {\it Utilitas Math}.

\bibitem{crc}
C.~J. Colbourn and J.~H. Dinitz.
\newblock {\em The CRC Handbook of combinatorial designs}.
\newblock CRC Press, 1996.

\bibitem{MR80j:05022}
D.~Curran and G.~H.~J. Van~Rees.
\newblock Critical sets in {L}atin squares.
\newblock In {\em Proceedings of the Eighth Manitoba Conference on Numerical
  Mathematics and Computing (Univ. Manitoba, Winnipeg, Man., 1978)}, pages
  165--168, Winnipeg, Man., 1979. Utilitas Math.

\bibitem{dk}
J.~Denes and D.~Keedwell.
\newblock {\em Latin squares and their applications}.
\newblock English Universities Press, 1974.

\bibitem{MR95k:05030}
D.~Donovan, J.~Cooper, D.~J. Nott, and J.~Seberry.
\newblock Latin squares: critical sets and their lower bounds.
\newblock {\em Ars Combin.}, {\bf 39} (1995), 33--48.

\bibitem{MR98b:05019}
D.~Donovan, A.~Howse, and P.~Adams.
\newblock A discussion of {L}atin interchanges.
\newblock {\em J. Combin. Math. Combin. Comput.}, {\bf 23} (1997), 161--182.

\bibitem{MR98m:05026}
A.~Howse.
\newblock Minimal critical sets for some small {L}atin squares.
\newblock {\em Australas. J. Combin.}, {\bf 17} (1998), 275--288.

\bibitem{MR99j:05036}
A.~D. Keedwell.
\newblock What is the size of the smallest {L}atin square for which a weakly
  completable critical set of cells exists?
\newblock {\em Ars Combin.}, {\bf 51} (1999), 97--104.

\bibitem{nauty}
B.~McKay.
\newblock {\em nauty} user's guide (version 1.5).
\newblock Australian National University Computer Science Technical Report
  TR-CS-90-02.

\end{thebibliography}

\end{document}
